% ==========================================================================
% 1. DOCUMENT SETUP AND PACKAGES
% ==========================================================================
% EDIT MAP
% Content: block 4, in page/section order; each row keeps its examples nearby.
% Typography/palettes/default spacing: block 2. All pt/mm values are LOGICAL
% A3 dimensions unless explicitly marked physical; printing applies 5/6 once.
% Highlights and measured table geometry: block 3. Preserve trailing % signs:
% they suppress real inter-cell glue, especially around overlay/measurement boxes.
% Local groups in block 4 deliberately override defaults; do not lift their
% settings into shared defaults. Native TikZ graphs remain beside their context.
\documentclass{article}
\usepackage[a3paper,portrait,hmargin=6mm,top=7mm,bottom=7mm]{geometry}
\usepackage{pgfpages}
\usepackage[T1]{fontenc}
\usepackage{lmodern}
\usepackage{amsmath,amssymb}
\usepackage{xcolor,colortbl}
\usepackage{tikz}
\usepackage{tabularx,array}
\usepackage{etoolbox} % Explicit dependency: \pretocmd patches table struts.
\usepackage{collcell} % Collect native math cells in composite example arrays.
% expl3 (\ExplSyntaxOn below) is provided by the current LaTeX kernel.

% ==========================================================================
% 2. SHARED SETTINGS: GEOMETRY, PALETTES, TYPOGRAPHY, SPACING
% ==========================================================================
% Retain the A3 composition as the logical typesetting canvas; the physical
% output below is exactly 250 x 350 mm. Fonts below account for the 5/6 scale.
% Uniform 5/6 scaling: 297 x 420 mm becomes 247.5 x 350 mm.
% Centering adds 1.25 mm at each side to absorb the aspect-ratio difference.
\pgfpagesdeclarelayout{poster-board-eighth}{}{%
  \pgfpagesphysicalpageoptions{logical pages=1,
    physical width=250mm,physical height=350mm}%
  \pgfpageslogicalpageoptions{1}{scale=0.833333333,
    center=\pgfpoint{125mm}{175mm}}%
}
\pgfpagesuselayout{poster-board-eighth}
\pagestyle{empty}
% Use scalable mathematical operators at the reference text sizes.
\input{omxlmex.fd}
\DeclareFontShape{OMX}{lmex}{m}{n}{<->lmex10}{}
% Central group palettes: dark ink, instructional accent, and quiet table rules.
% Group A retains the original colors; group B uses coordinated cool teals.
\definecolor{PosterAInk}{HTML}{172C3C}
\definecolor{PosterAAccent}{HTML}{BF604D}
\definecolor{PosterARule}{HTML}{A5AFB6}
\definecolor{PosterBInk}{HTML}{17383B}
\definecolor{PosterBAccent}{HTML}{287C7A}
\definecolor{PosterBRule}{HTML}{91AAA9}
% Existing semantic aliases also drive the unchanged panel tints and markers.
% Select after a page break so the preceding page retains its own palette.
\newcommand{\selectposterpalette}[1]{%
  \colorlet{PosterInk}{Poster#1Ink}%
  \colorlet{PosterAccent}{Poster#1Accent}%
  \colorlet{PosterRule}{Poster#1Rule}%
}
\selectposterpalette{A}
\color{PosterInk}
\setlength{\parindent}{0pt}
\setlength{\parskip}{0pt}
\setlength{\arrayrulewidth}{0.3pt}
% Central typography controls, explicitly applied in table cells.
% Printed sizes are 5/6 of these logical sizes: body 10.83 pt,
% examples 10.67 pt, symbols 11.67 pt, table labels/notes 10.33 pt.
\newcommand{\TitleFont}{\fontsize{27}{30}\selectfont\bfseries}
\newcommand{\KickerFont}{\fontsize{13}{15}\selectfont\bfseries}
\newcommand{\HeadingFont}{\fontsize{17}{20}\selectfont\bfseries}
\newcommand{\DefinitionFont}{\fontsize{13}{15}\selectfont}
\newcommand{\SymbolFont}{\fontsize{14}{16}\selectfont}
\newcommand{\ExampleFont}{\fontsize{12.8}{14.8}\selectfont}
\newcommand{\TableHeadFont}{\fontsize{12.4}{14.4}\selectfont\bfseries}
\newcommand{\FooterFont}{\fontsize{12.4}{14.4}\selectfont}
\DeclareMathSizes{12.6}{12.6}{8.8}{6.6}
\DeclareMathSizes{14}{14}{9.8}{7}
\DeclareMathSizes{13}{13}{9.1}{6.8}
\DeclareMathSizes{12.8}{12.8}{9}{6.8}


% Shared defaults: stretches are dimensionless; widths are linewidth fractions;
% RowPadding is extra clearance on EACH side, in addition to 0.4 logical mm.
\newcommand{\ExampleLineStretch}{1.12}
\newcommand{\TightExampleSpacing}{}
\newcommand{\MeasuredLineStretch}{1}
\newcommand{\ContextHeading}{Example contexts}
\newcommand{\SymbolColumnFraction}{.24}
\newcommand{\MeaningColumnFraction}{.32}
\newcommand{\PairExampleColumns}{r@{}E@{\quad}r@{}E}
\newlength{\RowPadding}
\setlength{\RowPadding}{0.7mm}
% Shared highlight switch; overlays never replace native text or mathematics.
\newif\ifPosterHighlights
\PosterHighlightstrue
% \PosterHighlightsfalse % Uncomment to disable instructional highlights.

% ==========================================================================
% 3. RENDERING HELPERS
% ==========================================================================
% 3.1 Native-metric highlight overlays (including integrals/fraction rules)
% One restrained marker for the specific notation introduced in each row.
% Rounded highlights keep even padding and the current math style.
% Native symbol metrics preserve spacing and baselines; padding uses surrounding clearance.
\newcommand{\notationhighlight}[1]{\ifPosterHighlights\mathpalette\highlightbackground{#1}\fi#1}
\newcommand{\highlightbackground}[2]{\begingroup
  % Paint only the background; the original symbol keeps its native math atom,
  % italic correction, script attachment, and dimensions outside this overlay.
  \raisebox{0pt}[0pt][0pt]{%
  \makebox[0pt][l]{\hspace{-1.2pt}\tikz[baseline=(notationhighlight.base)]\node[
    fill=PosterAccent!12,fill opacity=0.65,draw=none,rounded corners=2.6pt,
    inner xsep=1.2pt,inner ysep=2.2pt,outer sep=0pt,
    minimum width=0pt,minimum height=0pt
  ] (notationhighlight) {$\mathsurround=0pt #1\phantom{#2}$};}%
  }%
  \endgroup}
% Keep integral bounds on the native operator, including its italic correction;
% the shared renderer paints behind the integral sign only.
\newcommand{\notationintegral}[1]{\mathop{%
  \mathopen{\ifPosterHighlights\mathpalette\highlightbackground{\int}\fi}#1}\nolimits}
% Highlight the fraction rule alone, using the same renderer and toggle.
% Keep the native fraction (and its numerator/denominator spacing) intact.
\newsavebox{\FractionMeasure}
\newlength{\FractionBarWidth}
\newlength{\FractionAxis}
\newlength{\FractionRule}
\newcommand{\notationfraction}[2]{\mathinner{%
  \mathopen{\ifPosterHighlights\mathpalette\notationfractionbar{{#1}{#2}}\fi}%
  \frac{#1}{#2}}}
\newcommand{\notationfractionbar}[2]{\notationfractionmeasure{#1}#2}
\newcommand{\notationfractionmeasure}[3]{\begingroup
  \sbox{\FractionMeasure}{$\mathsurround=0pt #1\frac{#2}{#3}$}%
  \setlength{\FractionBarWidth}{\dimexpr\wd\FractionMeasure-2\nulldelimiterspace\relax}%
  \ifx#1\displaystyle
    \setlength{\FractionAxis}{\fontdimen22\textfont2}%
    \setlength{\FractionRule}{\fontdimen8\textfont3}%
  \else\ifx#1\textstyle
    \setlength{\FractionAxis}{\fontdimen22\textfont2}%
    \setlength{\FractionRule}{\fontdimen8\textfont3}%
  \else\ifx#1\scriptstyle
    \setlength{\FractionAxis}{\fontdimen22\scriptfont2}%
    \setlength{\FractionRule}{\fontdimen8\scriptfont3}%
  \else
    \setlength{\FractionAxis}{\fontdimen22\scriptscriptfont2}%
    \setlength{\FractionRule}{\fontdimen8\scriptscriptfont3}%
  \fi\fi\fi
  \raisebox{\dimexpr\FractionAxis-\FractionRule/2\relax}[0pt][0pt]{%
    \makebox[0pt][l]{\hspace{\nulldelimiterspace}%
      \highlightbackground{#1}{\rule{\FractionBarWidth}{\FractionRule}}}}%
  \endgroup}


% 3.2 Symmetric table struts (text tables only, never math arrays)
% Text tables use a symmetric minimum row strut around the cell centers.
% Keep its total height and leave mathematical arrays' internal spacing intact.
\makeatletter
\pretocmd{\@mkpream}{%
  \ifx\@acol\@tabacol
    \@tempdima=\dimexpr(\ht\@arstrutbox+\dp\@arstrutbox)/2\relax
    \setbox\@arstrutbox=\hbox{\vrule width0pt height\@tempdima depth\@tempdima}%
  \fi
  }{}{\PackageError{notation-posters}{Cannot center table struts}{}}
\makeatother

% 3.3 Measured cells, aligned examples, and table renderers
\newlength{\SymbolWidth}
\newlength{\MeaningWidth}
% Measure the complete cell (including wrapped text and tall mathematics),
% then reserve the same clearance above and below its actual bounds.
\newsavebox{\CellContent}
% All table helpers share the same vertical center and local padding.
% An extra 0.4 mm on each side gives A3 rows consistent breathing room.
% Dispatch hook: normally identical to centeredcell; calculus locally fixes
% header/reference clearance while body rows keep their own measured padding.
\newcommand{\paddedcell}[1]{\centeredcell{#1}}
\newcommand{\contextcell}[1]{%
  \sbox{\CellContent}{#1}%
  \raisebox{\dimexpr(\dp\CellContent-\ht\CellContent)/2\relax}{\usebox{\CellContent}}}
\newcommand{\centeredcell}[1]{%
  \sbox{\CellContent}{\begin{minipage}[t]{\linewidth}\raggedright #1\end{minipage}}%
  \raisebox{\dimexpr(\dp\CellContent-\ht\CellContent)/2\relax}%
    [\dimexpr(\ht\CellContent+\dp\CellContent)/2+\RowPadding+0.4mm\relax]%
    [\dimexpr(\ht\CellContent+\dp\CellContent)/2+\RowPadding+0.4mm\relax]%
    {\usebox{\CellContent}}%
}
% Keep each concept and its explanation in one measured Meaning cell.
% Paragraphs work in the cell minipages and the paired/triple parboxes alike.
\newcommand{\meaningexplanation}[1]{\textcolor{gray}{\textit{#1}}}
\newcommand{\conceptmeaning}[2]{\begingroup
  % Keep the concept and gray explanation close with readable A3 leading.
  % Inherit the current body leading, including the dense calculus page.
  #1\par #2\rule[-1pt]{0pt}{1pt}\par
  \endgroup}
\newsavebox{\ExampleContent}
% Text-style r/E arrays preserve relation spacing and native operator sizes.
\newcommand{\examplemath}[2][r@{}E]{\begingroup
  \renewcommand{\arraystretch}{\ExampleLineStretch}%
  \ifcsname NC@find@E\endcsname\else\newcolumntype{E}{>{{}}l}\fi%
  \sbox{\ExampleContent}{$\begin{array}[t]{@{}#1@{}}#2\end{array}$}%
  % Only overwide examples receive a small math-spacing correction.
  % This preserves the font size, baselines, and all fitting examples.
  \ifdim\wd\ExampleContent>\linewidth
    \thickmuskip=4mu\relax\medmuskip=3mu\relax
    \sbox{\ExampleContent}{$\begin{array}[t]{@{}#1@{}}#2\end{array}$}%
  \fi
  \ifdim\wd\ExampleContent>\linewidth
    \TightExampleSpacing
    \sbox{\ExampleContent}{$\begin{array}[t]{@{}#1@{}}#2\end{array}$}%
  \fi
  \usebox{\ExampleContent}\endgroup}
\newcommand{\notationrow}[3]{%
  \centeredcell{\centering\SymbolFont\ensuremath{\displaystyle #1}}&%
  \centeredcell{\DefinitionFont #2}&\centeredcell{\ExampleFont #3}\\\hline
}
% Composite rows retain one shared example array for horizontal relation alignment.
% Measure each logical row separately; center all its cells on that row's axis.
\newsavebox{\PairFirst}
\newsavebox{\PairSecond}
\newsavebox{\PairThird}
\newsavebox{\ExampleFirst}
\newsavebox{\ExampleSecond}
\newsavebox{\ExampleThird}
\newsavebox{\CurrentExample}
\newsavebox{\CurrentRow}
\newlength{\ExampleWidth}
\newcommand{\subrowcell}[1]{\begingroup
  \setlength{\RowPadding}{0pt}\centeredcell{#1}\endgroup}
\newcommand{\pairstrut}[1]{\rule[-\dp#1]{0pt}{\dimexpr\ht#1+\dp#1\relax}}
\newcommand{\measurelogicalrow}[6]{%
  \global\setbox#1=\hbox{\ExampleFont\examplemath[#3]{#6}}%
  \global\setbox#2=\hbox{%
    {\setlength{\linewidth}{\SymbolWidth}\subrowcell{\centering\SymbolFont$\displaystyle #4$}}%
    {\setlength{\linewidth}{\MeaningWidth}\subrowcell{\DefinitionFont #5}}%
    \raisebox{\dimexpr(\dp#1-\ht#1)/2\relax}{\usebox{#1}}}%
}
\newcommand{\pairlines}[3]{%
  \renewcommand{\arraystretch}{0}%
  \begin{tabular}[t]{@{}#1@{}}%
  \pairstrut{\PairFirst}\subrowcell{#2}\\%
  \pairstrut{\PairSecond}\subrowcell{#3}%
  \end{tabular}%
}
\newcommand{\triplelines}[4]{%
  \renewcommand{\arraystretch}{0}%
  \begin{tabular}[t]{@{}#1@{}}%
  \pairstrut{\PairFirst}\subrowcell{#2}\\%
  \pairstrut{\PairSecond}\subrowcell{#3}\\%
  \pairstrut{\PairThird}\subrowcell{#4}%
  \end{tabular}%
}
% Apply one translation to every math column in a logical row, preserving
% the relative baselines of its relations, fractions, and other mathematics.
\newcommand{\compositemathcell}[1]{%
  \pairstrut{\CurrentRow}%
  \raisebox{\dimexpr(\dp\CurrentExample-\ht\CurrentExample)/2\relax}{\renewcommand{\arraystretch}{\MeasuredLineStretch}$#1$}}
\newcommand{\compositerelationcell}[1]{\compositemathcell{{}#1}}
\newcolumntype{R}{>{\collectcell\compositemathcell}r<{\endcollectcell}}
\newcolumntype{L}{>{\collectcell\compositemathcell}l<{\endcollectcell}}
\newcolumntype{Q}{>{\collectcell\compositerelationcell}l<{\endcollectcell}}
\ExplSyntaxOn
\newcommand{\compositecolumns}[1]{
  \tl_set:Nn \l_tmpa_tl {#1}
  \tl_replace_all:Nnn \l_tmpa_tl {r}{R}
  \tl_replace_all:Nnn \l_tmpa_tl {l}{L}
  \tl_replace_all:Nnn \l_tmpa_tl {E}{Q}
  \cs_set:Npx \CompositeColumns {\exp_not:V \l_tmpa_tl}
}
\ExplSyntaxOff
\newcommand{\selectlogicalrow}[2]{%
  \global\setbox\CurrentRow=\copy#1%
  \global\setbox\CurrentExample=\copy#2}
\newcommand{\preparecomposite}{%
  \renewcommand{\ExampleLineStretch}{\MeasuredLineStretch}%
  \setlength{\ExampleWidth}{\dimexpr\linewidth-\SymbolWidth-\MeaningWidth-4\tabcolsep\relax}%
  \setlength{\linewidth}{\ExampleWidth}}
% Content interface: each braced entry contains {symbol}{meaning}{example}.
% This keeps one notation's content together. Columns are array preambles:
% r/E pairs share relation alignment; a single l preserves a nested array.
% Mandatory arguments let these commands expand directly to \noalign in tables.
\newcommand{\pairedentries}[2]{\expandafter\pairedentrieswithcolumns\expandafter{\PairExampleColumns}{#1}{#2}}
\newcommand{\pairedentrieswithcolumns}[3]{\unpackpairedentries{#1}#2#3}
\newcommand{\unpackpairedentries}[7]{%
  \renderpairedrow{#1}{#2}{#5}{#3}{#6}{#4}{#7}}
\newcommand{\tripleentries}[3]{\unpacktripleentries#1#2#3}
\newcommand{\unpacktripleentries}[9]{%
  \rendertriplerow{#1}{#4}{#7}{#2}{#5}{#8}{#3}{#6}{#9}}
% Internal renderers below accept column-major arguments solely to preserve
% the measured alignment algorithm. Edit entries in block 4, not these lists.
\newcommand{\rendertriplerow}[9]{%
  \noalign{\begingroup\preparecomposite
    \measurelogicalrow{\ExampleFirst}{\PairFirst}{r@{}E}{#1}{#4}{#7}%
    \measurelogicalrow{\ExampleSecond}{\PairSecond}{r@{}E}{#2}{#5}{#8}%
    \measurelogicalrow{\ExampleThird}{\PairThird}{r@{}E}{#3}{#6}{#9}%
  \endgroup}%
  \paddedcell{\centering\SymbolFont\triplelines{c}%
    {\centering$\displaystyle #1$}{\centering$\displaystyle #2$}{\centering$\displaystyle #3$}}&%
  \paddedcell{\DefinitionFont\triplelines{l}{#4}{#5}{#6}}&%
  \paddedcell{\ExampleFont\renewcommand{\ExampleLineStretch}{0}%
    \selectlogicalrow{\PairFirst}{\ExampleFirst}%
    \examplemath[R@{}Q]{\noalign{\selectlogicalrow{\PairFirst}{\ExampleFirst}}#7\\
      \noalign{\selectlogicalrow{\PairSecond}{\ExampleSecond}}#8\\
      \noalign{\selectlogicalrow{\PairThird}{\ExampleThird}}#9}}\\\hline
}
\newcommand{\renderpairedrow}[7]{%
  \noalign{\begingroup\preparecomposite
    \measurelogicalrow{\ExampleFirst}{\PairFirst}{#1}{#2}{#4}{#6}%
    \measurelogicalrow{\ExampleSecond}{\PairSecond}{#1}{#3}{#5}{#7}%
  \endgroup}%
  \paddedcell{\centering\SymbolFont\pairlines{c}{\centering$\displaystyle #2$}{\centering$\displaystyle #3$}}&%
  \paddedcell{\DefinitionFont\pairlines{l}{#4}{#5}}&%
  \paddedcell{\ExampleFont\renewcommand{\ExampleLineStretch}{0}%
    \compositecolumns{#1}%
    \def\SingleExampleColumn{l}\def\TheseExampleColumns{#1}%
    \ifx\SingleExampleColumn\TheseExampleColumns
      \def\FirstExample{{#6}}\def\SecondExample{{#7}}%
    \else\def\FirstExample{#6}\def\SecondExample{#7}\fi%
    \selectlogicalrow{\PairFirst}{\ExampleFirst}%
    \expandafter\examplemath\expandafter[\CompositeColumns]{\noalign{\selectlogicalrow{\PairFirst}{\ExampleFirst}}\FirstExample\\
      \noalign{\selectlogicalrow{\PairSecond}{\ExampleSecond}}\SecondExample}}\\\hline
}
% Compact contexts use the existing palette and the table's example type.
\newcommand{\examplecontexts}[1]{%
  \begingroup
  \setlength{\fboxsep}{0.5mm}%
  \setlength{\fboxrule}{0.35pt}%
  \noindent\fcolorbox{PosterAccent}{PosterAccent!10}{%
    \begin{minipage}{\dimexpr\linewidth-2\fboxsep-2\fboxrule\relax}
      \ExampleFont\color{PosterInk}%
      {\centering\textbf{\ContextHeading}\par}
      #1\par
    \end{minipage}}\par
  \endgroup
  \vspace{0.7mm}%
}
% A heading always spans the page; parallel tables belong to that section.
\newcommand{\sectionheading}[1]{%
  {\HeadingFont\color{PosterAccent}#1\par}%
  \vspace{0.4mm}{\color{PosterInk}\hrule height 0.5pt}\vspace{0.8mm}%
}
\newcommand{\notationtable}[1]{%
  \begingroup
  \setlength{\SymbolWidth}{\SymbolColumnFraction\linewidth}%
  \setlength{\MeaningWidth}{\MeaningColumnFraction\linewidth}%
  \setlength{\tabcolsep}{0.5mm}%
  \renewcommand{\arraystretch}{1}%
  \arrayrulecolor{PosterRule}%
  \begin{tabularx}{\linewidth}{@{}>{\centering\arraybackslash\SymbolFont}p{\SymbolWidth}>{\raggedright\arraybackslash\DefinitionFont}p{\MeaningWidth}>{\raggedright\arraybackslash\ExampleFont}X@{}}
  \paddedcell{\centering\TableHeadFont Symbol}&\paddedcell{\TableHeadFont Meaning}&\paddedcell{\TableHeadFont Examples}\\\hline
  #1
  \end{tabularx}%
  \endgroup
}
% Read the complete left table, then the right table, before the next section.
\newcommand{\paralleltables}[2]{%
  \noindent\begin{minipage}[t]{.49\linewidth}\vspace{0pt}\notationtable{#1}\end{minipage}\hfill
  \begin{minipage}[t]{.49\linewidth}\vspace{0pt}\notationtable{#2}\end{minipage}\par
  \vspace{1mm}%
}
% Parallel tables retain their measured subrows, local widths, and padding.
\newcommand{\compacttables}{%
  \renewcommand{\TightExampleSpacing}{\thinmuskip=0mu\relax\medmuskip=1mu\relax\thickmuskip=1mu\relax}%
  \setlength{\RowPadding}{0.9mm}%
  \renewcommand{\MeasuredLineStretch}{1}%
  \renewcommand{\SymbolColumnFraction}{.22}%
  \renewcommand{\MeaningColumnFraction}{.35}%
  \renewcommand{\PairExampleColumns}{r@{}E}%
  \renewcommand{\notationrow}[3]{%
    \paddedcell{\centering\SymbolFont\ensuremath{\displaystyle ##1}}&%
    \paddedcell{\DefinitionFont ##2}&\paddedcell{\ExampleFont ##3}\\\hline}%
}
% Each poster has its own fixed page and repeated identity.
\newcommand{\posterheader}[1]{%
  {\TitleFont\color{PosterAccent}MATHEMATICS\hfill\KickerFont #1\par}%
  \vspace{2mm}{\color{PosterInk}\hrule height 0.8pt}\vspace{3mm}%
  \color{PosterInk}%
}

% ==========================================================================
% 4. POSTER CONTENT: PDF READING ORDER
% ==========================================================================

% Five-column reference row: symbol, concept, explanation, three examples.
\newcommand{\comparisonrow}[6]{%
  \paddedcell{\centering\SymbolFont\ensuremath{\displaystyle #1}}&%
  \paddedcell{\DefinitionFont\conceptmeaning{#2}{\meaningexplanation{#3}}}&%
  \paddedcell{\ExampleFont #4}&\paddedcell{\ExampleFont #5}&\paddedcell{\ExampleFont #6}\\\hline}
\begin{document}
% ==========================================================================
% PAGE A-1: FOUNDATIONS
% ==========================================================================
\selectposterpalette{A}
\posterheader{A-1\quad FOUNDATIONS}
\begingroup
\begingroup
% General notation uses two independently headed tables in reading order.
\setlength{\RowPadding}{1.1mm}
\renewcommand{\MeasuredLineStretch}{1}
\renewcommand{\SymbolColumnFraction}{.26}
\renewcommand{\MeaningColumnFraction}{.33}
{\HeadingFont\color{PosterAccent}General notation\par}
\vspace{0.4mm}{\color{PosterInk}\hrule height 0.5pt}\vspace{0.4mm}
\noindent\begin{minipage}[t]{.49\linewidth}
% Balance the shorter left table using cell clearance only.
\setlength{\RowPadding}{0.9mm}
\vspace{0pt}
\notationtable{
\tripleentries%
  {% Logical entry: symbol, meaning, example
    {\mathrel{\notationhighlight{\mathord{=}}}}%
    {Equal}%
    {2+3&\mathrel{\notationhighlight{\mathord{=}}}5\quad 4-1\mathrel{\notationhighlight{\mathord{=}}}3}%
  }%
  {% Logical entry: symbol, meaning, example
    {\mathrel{\notationhighlight{\mathord{\neq}}}}%
    {Not equal}%
    {2&\mathrel{\notationhighlight{\mathord{\neq}}}3\quad 5+1\mathrel{\notationhighlight{\mathord{\neq}}}7}%
  }%
  {% Logical entry: symbol, meaning, example
    {\mathrel{\notationhighlight{\mathord{\approx}}}}%
    {Approximately}%
    {3.14&\mathrel{\notationhighlight{\mathord{\approx}}}\pi \quad\>\>\> \sqrt2\mathrel{\notationhighlight{\mathord{\approx}}}1.41}%
  }

\pairedentries%
  {% Logical entry: symbol, meaning, example
    {\mathrel{\notationhighlight{\mathord{<}}} \quad \mathrel{\notationhighlight{\mathord{>}}}}%
    {Less | Greater than}%
    {2&\mathrel{\notationhighlight{\mathord{<}}}3 & \qquad 5&\mathrel{\notationhighlight{\mathord{>}}}4}%
  }%
  {% Logical entry: symbol, meaning, example
    {\mathrel{\notationhighlight{\mathord{\nless}}} \quad \mathrel{\notationhighlight{\mathord{\ngtr}}}}%
    {Not Less | Greater}%
    {2&\mathrel{\notationhighlight{\mathord{\nless}}}2 &\qquad  4&\mathrel{\notationhighlight{\mathord{\ngtr}}}5}%
  }

\pairedentries%
  {% Logical entry: symbol, meaning, example
    {\mathrel{\notationhighlight{\mathord{\leq}}} \quad \mathrel{\notationhighlight{\mathord{\geq}}}}%
    {Less | Greater than or equal to}%
    {2&\mathrel{\notationhighlight{\mathord{\leq}}}2 &\qquad  5&\mathrel{\notationhighlight{\mathord{\geq}}}4}%
  }%
  {% Logical entry: symbol, meaning, example
    {\mathrel{\notationhighlight{\mathord{\nleq}}} \quad \mathrel{\notationhighlight{\mathord{\ngeq}}}}%
    {Not Less | Greater than or equal to}%
    {3&\mathrel{\notationhighlight{\mathord{\nleq}}}2 &\qquad  4&\mathrel{\notationhighlight{\mathord{\ngeq}}}5}%
  }

\notationrow
  {\begin{gathered}ab=a(b)\\{}=a\mathbin{\notationhighlight{\cdot}} b=a\mathbin{\notationhighlight{\times}} b\end{gathered}}
  {Multiplication: \\ \meaningexplanation{juxtaposition, parentheses, dot, or cross}}
  {\examplemath{3x&=3\mathbin{\notationhighlight{\cdot}} x=3\mathbin{\notationhighlight{\times}} x\\2(4)&=2\mathbin{\notationhighlight{\cdot}}4=2\mathbin{\notationhighlight{\times}}4}}
\notationrow
  {\begin{gathered}\notationfraction{a}{b}=a\notationhighlight{/}b\\{}=a\mathbin{\notationhighlight{\div}} b\end{gathered}}
  {Division: \\ \meaningexplanation{fraction, slash, or division sign ($b\neq0$)}}
  {\examplemath{\notationfraction{12}{3}&=12\notationhighlight{/}3=12\mathbin{\notationhighlight{\div}}3=4\\\notationfraction{x}{5}&=x\notationhighlight{/}5=x\mathbin{\notationhighlight{\div}}5\\& \qquad \quad =\notationfraction{1}{5}x \> \> \! =0.2x}}
}
\end{minipage}\hfill
\begin{minipage}[t]{.485\linewidth}
\vspace{0pt}
\notationtable{
\notationrow
  {\mathbin{\notationhighlight{\pm}}}
  {Plus or minus: \meaningexplanation{two choices}}
  {\examplemath{3\mathbin{\notationhighlight{\pm}}1&=4\text{ or }2\\-2\mathbin{\notationhighlight{\pm}}3&=1\text{ or }-5}}
\notationrow
  {\notationhighlight{|}a\notationhighlight{|}}
  {Absolute value: \meaningexplanation{Distance to $0$}}
  {\examplemath{\notationhighlight{|}-5\notationhighlight{|}&=5\\\notationhighlight{|}3\notationhighlight{|}&=3\quad \notationhighlight{|}0\notationhighlight{|}=0}}
\pairedentrieswithcolumns{l}%
  {% Logical entry: symbol, meaning, example
    {a^{\notationhighlight{2}}}%
    {Squared \par \meaningexplanation{Two factors of $a$}}%
    {\begin{array}[t]{@{}r@{}l@{}}3^{\notationhighlight{2}}&{}=3\cdot3=9\\(-2)^{\notationhighlight{2}}&{}=(-2)(-2)=4\end{array}}%
  }%
  {% Logical entry: symbol, meaning, example
    {a^{\notationhighlight{3}}}%
    {Cubed \par \meaningexplanation{Three factors of $a$}}%
    {\begin{array}[t]{@{}r@{}l@{}}2^{\notationhighlight{3}}&{}=2\cdot2\cdot2=8\\(-2)^{\notationhighlight{3}}&{}=(-2)(-2)(-2)\\&{}=-8\end{array}}%
  }
\pairedentrieswithcolumns{r@{}E}%
  {% Logical entry: symbol, meaning, example
    {\notationhighlight{\sqrt{a}}=a^{\notationhighlight{1/2}}}%
    {Square root \quad \> \meaningexplanation{$a\geq0$}}%
    {\notationhighlight{\sqrt9}&=3\text{ (not }\pm3\text{)}}%
  }%
  {% Logical entry: symbol, meaning, example
    {\sqrt[\notationhighlight{3}]{a}=a^{\notationhighlight{1/3}}}%
    {Cube root \qquad \! \meaningexplanation{$a\in\mathbb R$}}%
    {\sqrt[\notationhighlight{3}]{-8}&=-2}%
  }
\pairedentrieswithcolumns{l}%
  {% Logical entry: symbol, meaning, example
    {a^{\notationhighlight{n}}}%
    {$n$th power}%
    {\begin{array}[t]{@{}r@{}l@{}}2^{\notationhighlight{4}}&{}=2\cdot2\cdot2\cdot2=16\\(-2)^{\notationhighlight{5}}&{}=-32\end{array}}%
  }%
  {% Logical entry: symbol, meaning, example
    {\sqrt[\notationhighlight{n}]{a}=a^{1/\notationhighlight{n}}}%
    {$n$th root}%
    {\begin{array}[t]{@{}r@{}l@{}}\sqrt[\notationhighlight{4}]{16}&{}=\> \> \> 16^{\frac1{\notationhighlight{4}}} \quad  =2\\\sqrt[\notationhighlight{5}]{-32}&{}=(-32)^{\frac1{\notationhighlight{5}}}=-2\end{array}}%
  }
}
\end{minipage}\par\vspace{0mm}
\endgroup
\vspace{3mm}
\compacttables
\renewcommand{\MeaningColumnFraction}{.35}
% Shared heading row keeps both titles and their column rules aligned.
\noindent\begin{minipage}[t]{.49\linewidth}\vspace{0pt}
% --- \strut Sets ---
\sectionheading{\strut Sets}
\end{minipage}\hfill
\begin{minipage}[t]{.49\linewidth}\vspace{0pt}
% --- \strut Number systems ---
\sectionheading{\strut Number systems}
\end{minipage}\par\nobreak
\noindent\begin{minipage}[t]{.49\linewidth}\vspace{0pt}
\renewcommand{\ContextHeading}{Example guide}
\examplecontexts{$U$: Universal set \quad $U=\{1,2,3,4\}$ \> | \> $A=\{1,2\}$ \> | \> $B=\{2,3\}$}
% Fit the complement and its note using local cell clearance only.
\setlength{\RowPadding}{1mm}
\notationtable{
\notationrow
  {\mathopen{\notationhighlight{\{}}a,b,c\mathclose{\notationhighlight{\}}}}
  {\conceptmeaning{Set}{\meaningexplanation{Ignores order and repetition}}}
  {\examplemath{A&=\mathopen{\notationhighlight{\{}}2,1\mathclose{\notationhighlight{\}}}=\mathopen{\notationhighlight{\{}}1,2\mathclose{\notationhighlight{\}}}\\B&=\mathopen{\notationhighlight{\{}}2,3\mathclose{\notationhighlight{\}}}=\mathopen{\notationhighlight{\{}}2,2,3,3\mathclose{\notationhighlight{\}}}}}
\pairedentries%
  {% Logical entry: symbol, meaning, example
    {\mathrel{\notationhighlight{\mathord{\in}}}}%
    {Element of}%
    {2&\mathrel{\notationhighlight{\mathord{\in}}} A \qquad 2 \mathrel{\notationhighlight{\mathord{\in}}} \{2\}}%
  }%
  {% Logical entry: symbol, meaning, example
    {\mathrel{\notationhighlight{\mathord{\notin}}}}%
    {Not an element of}%
    {3&\mathrel{\notationhighlight{\mathord{\notin}}} A \qquad 3 \mathrel{\notationhighlight{\mathord{\notin}}} \{2\}}%
  }
\notationrow
  {\{x\mathrel{\notationhighlight{\mathord{:}}}p(x)\}}
  {\conceptmeaning{Set-builder notation}{\meaningexplanation{Select elements satisfying $p(x)$}}}
  {\examplemath[l]{A=\{x\in U\mathrel{\notationhighlight{\mathord{:}}}x\leq2\}\\B=\{x\in U\mathrel{\notationhighlight{\mathord{:}}}1<x<4\}\\\varnothing=\{x\in U\mathrel{\notationhighlight{\mathord{:}}}x>4\}}}
\pairedentries%
  {% Logical entry: symbol, meaning, example
    {\notationhighlight{\varnothing}}%
    {\conceptmeaning{Empty set}{\meaningexplanation{No elements}}}%
    {|\notationhighlight{\varnothing}|&=0}%
  }%
  {% Logical entry: symbol, meaning, example
    {\notationhighlight{|}A\notationhighlight{|}}%
    {\conceptmeaning{Cardinality}{\meaningexplanation{Number of elements}}}%
    {\notationhighlight{|}A\notationhighlight{|}&=2 \qquad \notationhighlight{|}U\notationhighlight{|}=4}%
  }
\pairedentrieswithcolumns{l}%
  {% Logical entry: symbol, meaning, example
    {A\mathrel{\notationhighlight{\mathord{\subseteq}}} U}%
    {\conceptmeaning{Subset}{\meaningexplanation{All elements included; equality allowed}}}%
    {\renewcommand{\arraystretch}{0.8}\begin{array}[t]{@{}l@{}}A\mathrel{\notationhighlight{\mathord{\subseteq}}} A \qquad \>\>\>\> A\mathrel{\notationhighlight{\mathord{\subseteq}}} U\\ \\ \{1\}\mathrel{\notationhighlight{\mathord{\subseteq}}}\{1,2\}\end{array}}%
  }%
  {% Logical entry: symbol, meaning, example
    {A\mathrel{\notationhighlight{\mathord{\subset}}} U}%
    {\conceptmeaning{Proper subset}{\meaningexplanation{Included but not equal}}}%
    {\renewcommand{\arraystretch}{0.8}\begin{array}[t]{@{}l@{}}\{1\}\mathrel{\notationhighlight{\mathord{\subset}}} A \qquad A\mathrel{\notationhighlight{\mathord{\subset}}} U\\ \\ A\mathrel{\notationhighlight{\mathord{\subset}}}\{1,2,3\}\end{array}}%
  }
\pairedentries%
  {% Logical entry: symbol, meaning, example
    {A\mathbin{\notationhighlight{\cup}} B}%
    {\conceptmeaning{Union}{\meaningexplanation{In either or both}}}%
    {A\mathbin{\notationhighlight{\cup}} B&=\{1,2,3\}}%
  }%
  {% Logical entry: symbol, meaning, example
    {A\mathbin{\notationhighlight{\cap}} B}%
    {\conceptmeaning{Intersection}{\meaningexplanation{In both}}}%
    {A\mathbin{\notationhighlight{\cap}} B&=\{2\}}%
  }
\notationrow
  {A\mathbin{\notationhighlight{\setminus}} B}
  {\conceptmeaning{Difference}{\meaningexplanation{In $A$, outside $B$}}}
  {\quad\>\>\! \examplemath{A & \mathbin{\notationhighlight{\setminus}} B=\{1\}} \\ \examplemath{\{3,4\} &\mathbin{\notationhighlight{\setminus}} \{4\}=\{3\}}}
\notationrow
  {A^{\notationhighlight{c}} \text{ or } A^{\notationhighlight{\prime}}}
  {\conceptmeaning{Complement}{\meaningexplanation{In $U$, outside $A$}}}
  {\examplemath{A^{\notationhighlight{c}}&=\{3,4\} \quad\>\>\>\>\>\> \{1\}^{\notationhighlight{c}} =  \{2,3,4\}\\B^{\notationhighlight{c}}&=\{1,4\} \quad\>\>\>\>\>\> \{1,2,3\}^{\notationhighlight{c}} =  \{4\}}}
}\par\vspace{1mm}
{\FooterFont \qquad\qquad $2$ is an element of $A$ \qquad | \qquad $\{2\}$ is a subset of $A$\par}
\end{minipage}\hfill
\begin{minipage}[t]{.49\linewidth}\vspace{0pt}
% Transfer three percent of this table from to Meaning.
\renewcommand{\SymbolColumnFraction}{.19}
\renewcommand{\MeaningColumnFraction}{.38}
% Fit the expanded number-system reference without changing typography.
\setlength{\RowPadding}{1mm}
\notationtable{
\pairedentries%
  {% Logical entry: symbol, meaning, example
    {\notationhighlight{\mathbb N}}%
    {Naturals}%
    {\notationhighlight{\mathbb N}&=\{1,2,3,\ldots\}}%
  }%
  {% Logical entry: symbol, meaning, example
    {\notationhighlight{\mathbb Z}}%
    {Integers}%
    {\notationhighlight{\mathbb Z}&=\{\ldots, -2, -1, 0, 1,2,\ldots\}}%
  }
\pairedentries%
  {% Logical entry: symbol, meaning, example
    {\notationhighlight{\mathbb Q}}%
    {\conceptmeaning{Rationals}{\meaningexplanation{$p/q$; $p,q\in\mathbb Z$, $q\ne0$}}}%
    {\tfrac12,-3&\in\notationhighlight{\mathbb Q}}%
  }%
  {% Logical entry: symbol, meaning, example
    {\notationhighlight{\mathbb R}}%
    {\conceptmeaning{Reals}{\meaningexplanation{The number line}}}%
    {\sqrt2,\pi,\tfrac12,-3&\in\notationhighlight{\mathbb R}}%
  }
\notationrow
  {\mathbb R\setminus\mathbb Q}
  {Irrationals}
  {\examplemath{\sqrt2,\pi&\in\mathbb R\setminus\mathbb Q}}
\notationrow
  {\notationhighlight{\mathbb C}}
  {\conceptmeaning{Complex numbers}{\meaningexplanation{$a+bi$; $a,b\in\mathbb R$, $i^2=-1$}}}
  {\examplemath{2+3i&\in\notationhighlight{\mathbb C}}}
\notationrow
  {\begin{gathered}[t]z=\\\notationhighlight{r}(\cos\notationhighlight{\theta}\\{}+i\sin\notationhighlight{\theta})\end{gathered}}
  {\conceptmeaning{Polar form:}{\meaningexplanation{$z=a+bi$:\\ $\notationhighlight{r}=|z|=\sqrt{a^2+b^2}$;\\ $\notationhighlight{\theta}$: argument in radians ($z\ne0$).}}}
  {\examplemath[l]{1+i=\notationhighlight{\sqrt{2}}(\cos(\notationhighlight{\pi/4})\\{}+i\sin(\notationhighlight{\pi/4}));\\\notationhighlight{r}=\notationhighlight{\sqrt2},\quad\notationhighlight{\theta}=\notationhighlight{\pi/4}}}
\pairedentries%
  {% Logical entry: symbol, meaning, example
    {\mathbb R^{\notationhighlight{2}}}%
    {\conceptmeaning{Real coordinate plane}{\meaningexplanation{Ordered pairs}}}%
    {(2,-3)&\in\mathbb R^{\notationhighlight{2}}}%
  }%
  {% Logical entry: symbol, meaning, example
    {\mathbb R^{\notationhighlight{3}}}%
    {\conceptmeaning{Real coordinate space}{\meaningexplanation{Ordered triples}}}%
    {(1,0,-2)&\in\mathbb R^{\notationhighlight{3}}}%
  }
\notationrow
  {\mathbb R^{\notationhighlight{+}}}
  {\conceptmeaning{Positive real numbers}{\meaningexplanation{Excludes zero}}}
  {\examplemath{\mathbb R^{\notationhighlight{+}}&=\{x\in\mathbb R:x>0\}\\&=(0,\infty)}}
\pairedentrieswithcolumns{l@{}E}%
  {% Logical entry: symbol, meaning, example
    {\mathopen{\notationhighlight{(}}a,b\mathclose{\notationhighlight{)}}}%
    {\conceptmeaning{Open interval}{\meaningexplanation{Excludes endpoints}}}%
    {x\in\mathopen{\notationhighlight{(}}1,3\mathclose{\notationhighlight{)}}:&\quad 1<x<3}%
  }%
  {% Logical entry: symbol, meaning, example
    {\mathopen{\notationhighlight{[}}a,b\mathclose{\notationhighlight{]}}}%
    {\conceptmeaning{Closed interval}{\meaningexplanation{Includes endpoints}}}%
    {x\in\mathopen{\notationhighlight{[}}1,3\mathclose{\notationhighlight{]}}:&\quad 1\leq x\leq3}%
  }
\pairedentrieswithcolumns{l}%
  {% Logical entry: symbol, meaning, example
    {\mathopen{\notationhighlight{[}}a,b\mathclose{\notationhighlight{)}}}%
    {\conceptmeaning{Half-open interval}{\meaningexplanation{Includes left endpoint}}}%
    {\begin{array}[t]{@{}l@{}}2\in\mathopen{\notationhighlight{[}}2,5\mathclose{\notationhighlight{)}} \qquad 5\notin\mathopen{\notationhighlight{[}}2,5\mathclose{\notationhighlight{)}}\end{array}}%
  }%
  {% Logical entry: symbol, meaning, example
    {\mathopen{\notationhighlight{(}}a,b\mathclose{\notationhighlight{]}}}%
    {\conceptmeaning{Half-open interval}{\meaningexplanation{Includes right endpoint}}}%
    {\begin{array}[t]{@{}l@{}}2\notin\mathopen{\notationhighlight{(}}2,5\mathclose{\notationhighlight{]}} \qquad 5\in\mathopen{\notationhighlight{(}}2,5\mathclose{\notationhighlight{]}}\end{array}}%
  }
\notationrow
  {-\notationhighlight{\infty}\;|\;\notationhighlight{\infty}}
  {Negative | Positive infinity}
  {\qquad \examplemath[l]{(-\notationhighlight{\infty},0)\quad|\quad[0,\notationhighlight{\infty})}}
}\vspace{1mm}
{\FooterFont Intervals assume $a<b$. Some texts include $0$ in $\mathbb N$\par}
\end{minipage}\par
\par\vspace{3mm}
\compacttables
% --- Greek symbols and constants ---
\sectionheading{Greek symbols and constants}
% One alphabetic letter table beside a separate constants table.
\begingroup
\setlength{\RowPadding}{0.15mm}
\setlength{\tabcolsep}{0.7mm}
\renewcommand{\arraystretch}{0.8}
\arrayrulecolor{PosterRule}
\newcommand{\greekentry}[3]{%
  \paddedcell{\DefinitionFont #1}&%
  \paddedcell{\centering\SymbolFont\ensuremath{#2}}&%
  \paddedcell{\centering\SymbolFont\ensuremath{#3}}%
}
\newcommand{\greekhead}{%
  \paddedcell{\TableHeadFont Name}&%
  \paddedcell{\centering\TableHeadFont Lowercase Symbol}&%
  \paddedcell{\centering\TableHeadFont Uppercase Symbol}%
}
\noindent\begin{minipage}[t]{.49\linewidth}
\vspace{0pt}
\setlength{\tabcolsep}{0.25mm}
\setlength{\RowPadding}{0.9mm}
% Read down the left group, then the right group: six letters each.
\begin{tabularx}{\linewidth}{@{}>{\raggedright\arraybackslash}X>{\centering\arraybackslash}p{.176\linewidth}>{\centering\arraybackslash}p{.176\linewidth}@{\hspace{0.7mm}{\color{lightgray}\vrule width 1.5pt}\hspace{0.7mm}}>{\raggedright\arraybackslash}X>{\centering\arraybackslash}p{.176\linewidth}>{\centering\arraybackslash}p{.176\linewidth}@{}}
\greekhead&\greekhead\\\hline
\greekentry{Alpha}{\alpha}{\mathrm A}&\greekentry{Mu}{\mu}{\mathrm M}\\\hline
\greekentry{Beta}{\beta}{\mathrm B}&\greekentry{Pi}{\pi}{\Pi}\\\hline
\greekentry{Gamma}{\gamma}{\Gamma}&\greekentry{Rho}{\rho}{\mathrm P}\\\hline
\greekentry{Delta}{\delta}{\Delta}&\greekentry{Sigma}{\sigma}{\Sigma}\\\hline
\greekentry{Theta}{\theta}{\Theta}&\greekentry{Phi}{\phi}{\Phi}\\\hline
\greekentry{Lambda}{\lambda}{\Lambda}&\greekentry{Omega}{\omega}{\Omega}\\\hline
\end{tabularx}
\end{minipage}\hfill
\begin{minipage}[t]{.49\linewidth}
\vspace{0pt}
\setlength{\RowPadding}{1.1mm}
\begin{tabularx}{\linewidth}{@{}>{\centering\arraybackslash}p{.16\linewidth}>{\raggedright\arraybackslash}p{.26\linewidth}>{\raggedright\arraybackslash}X@{}}
\paddedcell{\centering\TableHeadFont Constant}&
\paddedcell{\TableHeadFont Name}&
\paddedcell{\TableHeadFont Meaning / Value}\\\hline
\paddedcell{\centering\SymbolFont $\notationhighlight{\pi}$}&
\paddedcell{\DefinitionFont Pi}&
\paddedcell{\DefinitionFont Circle constant\par $\notationhighlight{\pi}=C/d\approx\notationhighlight{3.14159}$\par \meaningexplanation{$C$: circumference; $d$: diameter.}}\\\hline
\paddedcell{\centering\SymbolFont $\notationhighlight{e}$}&
\paddedcell{\DefinitionFont Euler's number}&
\paddedcell{\DefinitionFont Base of natural logarithms\par $\notationhighlight{e}\approx\notationhighlight{2.71828}$\par \meaningexplanation{$\ln x=y\iff \notationhighlight{e}^{y}=x$ ($x>0$);\\ $\ln1=0$, $\ln\notationhighlight{e}=1$}}\\\hline
\paddedcell{\centering\SymbolFont $\notationhighlight{i}$}&
\paddedcell{\DefinitionFont Imaginary unit}&
\paddedcell{\DefinitionFont $\notationhighlight{i}^2=-1$\par \meaningexplanation{$a+b\notationhighlight{i}$; $a,b\in\mathbb R$}}\\\hline
\end{tabularx}\par\vspace{0.8mm}
\end{minipage}
\endgroup
{\color{PosterInk}\hrule height 0.5pt}\vspace{0.7mm}
\endgroup
\newpage
% ==========================================================================
% PAGE A-2: FUNCTIONS AND PROBABILITY
% ==========================================================================
\selectposterpalette{A}
\posterheader{A-2\quad FUNCTIONS AND PROBABILITY}
\begingroup
\compacttables
\begingroup
% Compact spacing and balanced columns apply only to this middle poster.
\setlength{\RowPadding}{0.8mm}
\renewcommand{\SymbolColumnFraction}{.24}
\renewcommand{\MeaningColumnFraction}{.38}
% --- Function notation ---
\sectionheading{Function notation}
\begingroup
\setlength{\tabcolsep}{0.5mm}
\setlength{\RowPadding}{0.1mm}
\renewcommand{\arraystretch}{1}
\renewcommand{\ExampleLineStretch}{1.02}
\arrayrulecolor{PosterRule}
% Separate notation and definitions beside the three original example columns.

\begin{tabularx}{\linewidth}{@{}p{.14\linewidth}p{.22\linewidth}p{.20\linewidth}p{.20\linewidth}X@{}}
\paddedcell{\centering\TableHeadFont Symbol}&
\paddedcell{\TableHeadFont Meaning}&
\paddedcell{\TableHeadFont Square function}&
\paddedcell{\TableHeadFont Square-root function}&
\paddedcell{\TableHeadFont Linear function}\\\hline
\multicolumn{2}{@{}p{\dimexpr.36\linewidth+2\tabcolsep\relax}}{%
\paddedcell{\ExampleFont $h$: function; $x$: input; $h(x)$: output.\par
$(x,h(x))$ is a point on its graph.\par
A function assigns exactly one output in $K$ to each input in $D$.\par
\textbf{Piecewise function:} different formulas for different input conditions.\par
\noindent\begin{minipage}[c]{.56\linewidth}
$\displaystyle p(x)=\begin{cases}
  -x, & \notationhighlight{x<0},\\
  x^2, & \notationhighlight{x\ge0}.
\end{cases}$
\end{minipage}\hfill
\begin{minipage}[c]{.42\linewidth}
$p(-2)=\notationhighlight{2}$\par
$p(0)=\notationhighlight{0}$\par
$p(2)=\notationhighlight{4}$
\end{minipage}\par}&
% Native TikZ coordinates retain the numerical values from Function value below.
% Labels remain 9 pt at the physical 5/6 page size; only local graph units grow.
% Quiet projections are drawn first so curves and marked points stay prominent.
\paddedcell{\ExampleFont $f(x)=x^2$\par
  \begin{tikzpicture}[font=\fontsize{10.8}{12}\selectfont,
    every node/.style={inner sep=0pt},line cap=round,
    guide/.style={draw=PosterInk!45,line width=0.50pt,dash pattern=on 1.2pt off 1.2pt},
    leader/.style={draw=PosterInk!65,line width=0.6pt}]
    \path[use as bounding box] (0,0) rectangle (56mm,44mm);
    \begin{scope}[shift={(29mm,7mm)},x=5mm,y=2.2mm]
      \draw[guide] (-2,0)--(-2,4)--(0,4);
      \draw[guide] (3,0)--(3,9)--(0,9);
      \draw[->,line width=0.70pt] (-3.4,0)--(4.6,0);
      \draw[->,line width=0.70pt] (0,-.5)--(0,11);
      \draw[PosterAccent,line width=1.4pt,domain=-3.1:3.4,samples=65]
        plot (\x,{\x*\x});
      \fill (-2,4) circle[radius=1.1pt] (3,9) circle[radius=1.1pt];
      \node[below=1.5mm] at (-2,0) {$-2$};
      \node[below=1.5mm] at (3,0) {$3$};
      \node[right=1.5mm] at (0,4) {$4$};
      \node[left=1.5mm] at (0,9) {$9$};
      \coordinate (fNegativePoint) at (-2,4);
      \coordinate (fPositivePoint) at (3,9);
    \end{scope}
    \node[anchor=center] (fNegativeLabel) at (19mm,37mm) {$(-2,f(-2))$};
    \node[anchor=center] (fPositiveLabel) at (37.5mm,33mm) {$(3,f(3))$};
    \draw[leader,shorten <=1.5pt,shorten >=1mm]
      (fNegativePoint)--(fNegativeLabel);
    \draw[leader,shorten <=1.5pt,shorten >=1mm]
      (fPositivePoint)--(fPositiveLabel);
  \end{tikzpicture}}&
\paddedcell{\ExampleFont $g(x)=\sqrt{x}$\par
  \begin{tikzpicture}[font=\fontsize{10.8}{12}\selectfont,
    every node/.style={inner sep=0pt},line cap=round,
    guide/.style={draw=PosterInk!45,line width=0.50pt,dash pattern=on 1.2pt off 1.2pt},
    leader/.style={draw=PosterInk!65,line width=0.6pt}]
    \path[use as bounding box] (0,0) rectangle (56mm,44mm);
    \begin{scope}[shift={(6mm,7mm)},x=2.5mm,y=5.5mm]
      \draw[guide] (9,0)--(9,3)--(0,3);
      \draw[guide] (16,0)--(16,4)--(0,4);
      \draw[->,line width=0.70pt] (-.5,0)--(19,0);
      \draw[->,line width=0.70pt] (0,-.2)--(0,5.2);
      % Parameterization x=t^2, y=t gives a smooth start, only for x>=0.
      \draw[PosterAccent,line width=1.4pt,domain=0:4.25,samples=65,variable=\t]
        plot ({\t*\t},\t);
      \fill (9,3) circle[radius=1.1pt] (16,4) circle[radius=1.1pt];
      \node[below=1.5mm] at (9,0) {$9$};
      \node[below=1.5mm] at (16,0) {$16$};
      \node[left=1.5mm] at (0,3) {$3$};
      \node[left=1.5mm] at (0,4) {$4$};
      \coordinate (gNinePoint) at (9,3);
      \coordinate (gSixteenPoint) at (16,4);
    \end{scope}
    \node[anchor=center] (gNineLabel) at (37.25mm,17mm)
      {$(9,g(9))$};
    \node[anchor=center] (gSixteenLabel) at (35mm,35mm)
      {$(16,g(16))$};
    \draw[leader,shorten <=1.5pt,shorten >=1mm]
      (gNinePoint)--(gNineLabel);
    \draw[leader,shorten <=1.5pt,shorten >=1mm]
      (gSixteenPoint)--(gSixteenLabel);
  \end{tikzpicture}}&
\paddedcell{\ExampleFont $\beta(x)=2x+1$\par
  \begin{tikzpicture}[font=\fontsize{10.8}{12}\selectfont,
    every node/.style={inner sep=0pt},line cap=round,
    guide/.style={draw=PosterInk!45,line width=0.50pt,dash pattern=on 1.2pt off 1.2pt},
    leader/.style={draw=PosterInk!65,line width=0.6pt}]
    \path[use as bounding box] (0,0) rectangle (56mm,44mm);
    \begin{scope}[shift={(25mm,16mm)},x=6mm,y=2.7mm]
      \draw[guide] (-1,0)--(-1,-1)--(0,-1);
      \draw[guide] (3,0)--(3,7)--(0,7);
      \draw[->,line width=0.70pt] (-3.6,0)--(4.2,0);
      \draw[->,line width=0.70pt] (0,-2.2)--(0,8.5);
      \draw[PosterAccent,line width=1.4pt,domain=-1.6:3.7]
        plot (\x,{2*\x+1});
      \fill (-1,-1) circle[radius=1.1pt] (3,7) circle[radius=1.1pt];
      \node[above=1.5mm] at (-1,0) {$-1$};
      \node[below=1.5mm] at (3,0) {$3$};
      \node[right=1.5mm] at (0,-1) {$-1$};
      \node[left=1.5mm] at (0,7) {$7$};
      \coordinate (betaNegativePoint) at (-1,-1);
      \coordinate (betaPositivePoint) at (3,7);
    \end{scope}
    % Convert the former corner coordinates to centers using the label dimensions.
    \setbox0=\hbox{\fontsize{10.8}{12}\selectfont $(-1,\beta(-1))$}%
    \node[anchor=center] (betaNegativeLabel) at
      (\dimexpr23mm-.5\wd0-.5\pgflinewidth\relax,
       \dimexpr29mm-.5\ht0-.5\dp0-.5\pgflinewidth\relax)
      {$(-1,\beta(-1))$};
    \setbox0=\hbox{\fontsize{10.8}{12}\selectfont $(3,\beta(3))$}%
    \node[anchor=center] (betaPositiveLabel) at
      (\dimexpr28mm+.5\wd0+.5\pgflinewidth\relax,
       \dimexpr38mm+.5\ht0+.5\dp0+.5\pgflinewidth\relax)
      {$(3,\beta(3))$};
    \draw[leader,shorten <=1mm,shorten >=1.5pt]
      (betaNegativeLabel.south)--(betaNegativeLabel.south |- 0,11mm)--(betaNegativePoint);
    \draw[leader,shorten <=1.5pt,shorten >=1mm]
      (betaPositivePoint)--(betaPositiveLabel);
  \end{tikzpicture}}\\\hline
\comparisonrow
  {h:D\mathrel{\notationhighlight{\mathord{\to}}} K}
  {Function mapping}
  {Domain $D$ \quad codomain $K$}
  {$f:\mathbb R\mathrel{\notationhighlight{\mathord{\to}}}\mathbb R$}
  {$g:[0,\infty)\mathrel{\notationhighlight{\mathord{\to}}}\mathbb R$}
  {$\beta:\mathbb R\mathrel{\notationhighlight{\mathord{\to}}}\mathbb R$}
\comparisonrow
  {\mathop{\notationhighlight{\mathord{\operatorname{dom}}}}(h)=\notationhighlight{D}}
  {Domain}
  {Allowed inputs}
  {$\mathop{\notationhighlight{\mathord{\operatorname{dom}}}}(f)=\notationhighlight{\mathbb R}$}
  {$\mathop{\notationhighlight{\mathord{\operatorname{dom}}}}(g)=\notationhighlight{[0,\infty)}$}
  {$\mathop{\notationhighlight{\mathord{\operatorname{dom}}}}(\beta)=\notationhighlight{\mathbb R}$}
\comparisonrow
  {\mathop{\notationhighlight{\mathord{\operatorname{codom}}}}(h)=\notationhighlight{K}}
  {Codomain}
  {Declared output target}
  {$\mathop{\notationhighlight{\mathord{\operatorname{codom}}}}(f)=\notationhighlight{\mathbb R}$}
  {$\mathop{\notationhighlight{\mathord{\operatorname{codom}}}}(g)=\notationhighlight{\mathbb R}$}
  {$\mathop{\notationhighlight{\mathord{\operatorname{codom}}}}(\beta)=\notationhighlight{\mathbb R}$}
\comparisonrow
  {x\mathrel{\notationhighlight{\mathord{\mapsto}}} h(x)}
  {Mapping rule}
  {$\mapsto$ maps inputs\enspace $\to$ links sets}
  {\quad\>\!\! $x\mathrel{\notationhighlight{\mathord{\mapsto}}} x^2$}
  {\quad\>\!\! $x\mathrel{\notationhighlight{\mathord{\mapsto}}}\sqrt{x}$}
  {\quad\>\!\! $x\mathrel{\notationhighlight{\mathord{\mapsto}}} 2x+1$}
\comparisonrow
  {\notationhighlight{h(x)}}
  {Function value}
  {Output at input $x$}
  {\examplemath[l]{f\mathopen{\notationhighlight{(}}3\mathclose{\notationhighlight{)}}=3^2=\notationhighlight{9}\\f\mathopen{\notationhighlight{(}}-2\mathclose{\notationhighlight{)}}=(-2)^2=\notationhighlight{4}}}
  {\examplemath[l]{g\mathopen{\notationhighlight{(}}9\mathclose{\notationhighlight{)}}=\sqrt{9}=\notationhighlight{3}\\g\mathopen{\notationhighlight{(}}16\mathclose{\notationhighlight{)}}=\sqrt{16}=\notationhighlight{4}}}
  {\examplemath[l]{\beta\mathopen{\notationhighlight{(}}3\mathclose{\notationhighlight{)}}=2(3)+1=\notationhighlight{7}\\\beta\mathopen{\notationhighlight{(}}-1\mathclose{\notationhighlight{)}}=2(-1)+1=\notationhighlight{-1}}}
\comparisonrow
  {\mathop{\notationhighlight{\mathord{\operatorname{range}}}}(h)}
  {Range}
  {Only attained outputs}
  {$\mathop{\notationhighlight{\mathord{\operatorname{range}}}}(f)=\notationhighlight{[0,\infty)}$}
  {$\mathop{\notationhighlight{\mathord{\operatorname{range}}}}(g)=\notationhighlight{[0,\infty)}$}
  {$\mathop{\notationhighlight{\mathord{\operatorname{range}}}}(\beta)=\notationhighlight{\mathbb R}$}
\end{tabularx}\par
\endgroup
\vspace{2mm}
% --- Probability and events ---
\sectionheading{Probability and events}
\begingroup
\setlength{\tabcolsep}{0.6mm}
% Share exact column dimensions between the context panel and main table.
\setlength{\SymbolWidth}{.13\linewidth}
\setlength{\MeaningWidth}{.25\linewidth}
\setlength{\ExampleWidth}{\dimexpr(\linewidth-\SymbolWidth-\MeaningWidth-6\tabcolsep)/2\relax}
\newcolumntype{V}{@{}p{\SymbolWidth}p{\MeaningWidth}p{\ExampleWidth}p{\ExampleWidth}@{}}
\begingroup
\setlength{\fboxsep}{0.5mm}
\setlength{\fboxrule}{0.35pt}
\noindent\fcolorbox{PosterAccent}{PosterAccent!10}{%
  \begin{minipage}{\dimexpr\linewidth-2\fboxsep-2\fboxrule\relax}
  \ExampleFont\color{PosterInk}%
  % Cancel the frame inset to retain the main table's column origins.
  \makebox[\linewidth][l]{%
  \hspace*{\dimexpr-\fboxsep-\fboxrule\relax}%
  \begin{tabular}{V}
  \multicolumn{2}{@{}p{\dimexpr\SymbolWidth+\MeaningWidth+2\tabcolsep\relax}}{%
    \begin{minipage}[t]{\linewidth}\raggedright
    {\centering\textbf{Example contexts}\par}
    \vspace{0.5mm}
    \textbf{Fractions, decimals, and percentages can express the same probability:}\par
    \vspace{0.5mm}
    {\centering
    $\displaystyle P(A)=\frac{1}{2}=0.5=50\%$\par
    \vspace{0.5mm}
    $\displaystyle P(F)=\frac{3}{4}=0.75=75\%$\par}
    \vspace{0.5mm}
    Divide the numerator by the denominator to obtain the decimal representation.
    Multiply the decimal by $100$ to obtain the percentage.\par
    \end{minipage}}&
  \begin{minipage}[t]{\ExampleWidth}\raggedright
  \textbf{Die: one fair die roll}\par
  A six-sided die\par
  $\Omega=\{1,2,3,4,5,6\}$\par
  \begingroup
  \begin{tabularx}{\linewidth}{@{}>{\collectcell\contextcell}l<{\endcollectcell}>{\collectcell\subrowcell}X<{\endcollectcell}@{}}
  \textbf{Symbol / event set}&\textbf{Meaning}\\
  $A=\{1,3,5\}$&Odd result\\
  $B=\{1,2\}$&Result at most $2$\\
  $C=\{2,4,6\}$&Even result
  \end{tabularx}\endgroup
  \end{minipage}&
  \begin{minipage}[t]{\ExampleWidth}\raggedright
  \textbf{Coins: two independent fair tosses}\par
  $\mathrm H$: heads, $\mathrm T$: tails\par
  $\Omega=\{\mathrm{HH},\mathrm{HT},\mathrm{TH},\mathrm{TT}\}$\par
  \begingroup
  \begin{tabularx}{\linewidth}{@{}>{\collectcell\contextcell}l<{\endcollectcell}>{\collectcell\subrowcell}X<{\endcollectcell}@{}}
  \textbf{Symbol / event set}&\textbf{Meaning}\\
  $F=\{\mathrm{HH},\mathrm{HT},\mathrm{TH}\}$&At least one head\\
  $G=\{\mathrm{HH},\mathrm{TT}\}$&Matching results\\
  $J=\{\mathrm{TT}\}$&No heads
  \end{tabularx}\endgroup
  \end{minipage}
  \end{tabular}}\par
  \end{minipage}}\par
\endgroup
\vspace{0.7mm}
% Compact this table locally so the complete probability poster stays on one page.
\setlength{\RowPadding}{0.15mm}
\renewcommand{\arraystretch}{1}
\renewcommand{\ExampleLineStretch}{1.02}
\arrayrulecolor{PosterRule}
% Separate notation and definitions beside the two original experiment columns.
\newcommand{\probabilityrow}[5]{%
  \paddedcell{\centering\SymbolFont\ensuremath{\displaystyle #1}}&%
  \paddedcell{\DefinitionFont\conceptmeaning{\if\relax\detokenize{#2}\relax\else#2:\fi}{\meaningexplanation{#3}}}&%
  \paddedcell{\ExampleFont #4}&\paddedcell{\ExampleFont #5}\\\hline}
\begin{tabular}{V}
\paddedcell{\centering\TableHeadFont Symbol}&
\paddedcell{\TableHeadFont Meaning}&
\paddedcell{\TableHeadFont Die examples}&
\paddedcell{\TableHeadFont Coins examples}\\\hline
\probabilityrow
  {\notationhighlight{\Omega}\quad\notationhighlight{\omega}}
  {Sample space; outcome}
  {All outcomes; one outcome}
  {$|\notationhighlight{\Omega}|=6$ \qquad $\notationhighlight{\omega}=3\in\notationhighlight{\Omega}$\par \qquad \qquad \quad \> $\notationhighlight{\omega}=6\in\notationhighlight{\Omega}$}
  {$|\notationhighlight{\Omega}|=4$ \qquad $\notationhighlight{\omega}=\mathrm{HT}\in\notationhighlight{\Omega}$\par \qquad \qquad \quad \> $\notationhighlight{\omega}=\mathrm{HH}\in\notationhighlight{\Omega}$}
\probabilityrow
  {\notationhighlight{E}\subseteq\Omega}
  {Input: event}
  {Outcome set, not a probability}
  {\examplemath[l]{\notationhighlight{E}_1=\{1,4,6\},\quad |\notationhighlight{E}_1|=3\\\notationhighlight{E}_2=\{2,5\},\quad |\notationhighlight{E}_2|=2}}
  {\examplemath[l]{\notationhighlight{E}_1=\{\mathrm{HH},\mathrm{HT}\},\quad |\notationhighlight{E}_1|=2\\\notationhighlight{E}_2=\{\mathrm{HH},\mathrm{TH}\},\quad |\notationhighlight{E}_2|=2}}
\probabilityrow
  {\begin{gathered}\notationhighlight{\mathcal P}(\Omega)\\{}=\{E:E\subseteq\Omega\}\end{gathered}}
  {Power set; domain of $P$}
  {All events (subsets of $\Omega$).}
  {$|\notationhighlight{\mathcal P}(\Omega)|=2^6=64$; e.g.\par $\varnothing,\Omega,A,B,C,\{1\},A\cap B$}
  {$|\notationhighlight{\mathcal P}(\Omega)|=2^4=16$; e.g.\par $\varnothing,\Omega,F,G,J,\{\mathrm{HT}\},F\cap G$}
\probabilityrow
  {\notationhighlight{P}}
  {Probability measure}
  {Function $\notationhighlight{P}:\mathcal P(\Omega)\to[0,1]$, $E\mapsto \notationhighlight{P}(E)$}
  {\examplemath[l@{\quad}l]{\notationhighlight{P}(\varnothing)=0&\notationhighlight{P}(A)=\tfrac12\\\notationhighlight{P}(B)=\tfrac13&\notationhighlight{P}(\{6\})=\tfrac16}}
  {\examplemath[l@{\quad}l]{\notationhighlight{P}(\Omega)=1&\notationhighlight{P}(F)=\tfrac34\\\notationhighlight{P}(G)=\tfrac12&\notationhighlight{P}(\{\mathrm{HT}\})=\tfrac14}}
\probabilityrow
  {[0,1]}
  {Codomain}
  {Declared numerical target }
  {$[0,1]$ or $[0\%,100\%]$}
  {$[0,1]$ $[0,1]$ or $[0\%,100\%]$}

\probabilityrow
  {\notationhighlight{P(E)}}
  {Output: probability}
  {Numerical value; $P(\varnothing)=\notationhighlight{0}$}
  {$P(A)=\notationhighlight{\tfrac12}$ \quad $P(B)=\notationhighlight{\tfrac13}$ \quad $P(C)=\notationhighlight{\tfrac12}$}
  {$P(F)=\notationhighlight{\tfrac34}$ \quad $P(G)=\notationhighlight{\tfrac12}$ \quad $P(J)=\notationhighlight{\tfrac14}$}
\probabilityrow
  {\notationhighlight{\text{Range}}(P)}
  {}
  {Attained probabilities}
  {$\notationhighlight{\{0,\tfrac16,\tfrac26,\ldots,1\}}$}
  {$\notationhighlight{\{0,\tfrac14,\tfrac12,\tfrac34,1\}}$}
\probabilityrow
  {P(E)=\frac{\notationhighlight{|}E\notationhighlight{|}}{\notationhighlight{|}\Omega\notationhighlight{|}}}
  {Classical probability}
  {Equally likely finite outcomes:}
  {\examplemath[l]{A=\{1,3,5\}\\P(A)=\frac{\notationhighlight{|}A\notationhighlight{|}}{\notationhighlight{|}\Omega\notationhighlight{|}}=\frac{\notationhighlight{3}}{\notationhighlight{6}}=\tfrac12}}
  {\examplemath[l]{F=\{\mathrm{HH},\mathrm{HT},\mathrm{TH}\}\\P(F)=\frac{\notationhighlight{|}F\notationhighlight{|}}{\notationhighlight{|}\Omega\notationhighlight{|}}=\frac{\notationhighlight{3}}{\notationhighlight{4}}}}
\probabilityrow
  {E^{\notationhighlight{c}}}
  {Complement}
  {In $\Omega$, outside $E$; $P(E^{\notationhighlight{c}})=1-P(E)$}
  {\examplemath[l@{\quad}l]{A^{\notationhighlight{c}}=C&\{1,2\}^{\notationhighlight{c}}=\{3,4,5,6\}\\P(A^{\notationhighlight{c}})=\tfrac12&P(\{1,2\}^{\notationhighlight{c}})=\tfrac46=\tfrac23}}
  {\examplemath[l@{\quad}l]{F^{\notationhighlight{c}}=J&\{\mathrm{HT}\}^{\notationhighlight{c}}=\{\mathrm{HH},\mathrm{TH},\mathrm{TT}\}\\P(F^{\notationhighlight{c}})=\tfrac14&P(\{\mathrm{HT}\}^{\notationhighlight{c}})=\tfrac34}}
\probabilityrow
  {E_1\mathbin{\notationhighlight{\cup}} E_2}
  {Union}
  {$E_1$ or $E_2$, or both}
  {$A\mathbin{\notationhighlight{\cup}} B=A+B=\{1,2,3,5\}$}
  {$F\mathbin{\notationhighlight{\cup}} G=F+G=\Omega$}
\probabilityrow
  {E_1\mathbin{\notationhighlight{\cap}} E_2}
  {Intersection}
  {$E_1$ and $E_2$ together}
  {$A\mathbin{\notationhighlight{\cap}} B=AB=\{1\}$}
  {$F\mathbin{\notationhighlight{\cap}} G=FG=\{\mathrm{HH}\}$}
\probabilityrow
  {\notationhighlight{E_1\cap E_2=\varnothing}}
  {Disjoint}
  {Mutually exclusive:}
  {$\notationhighlight{A\cap C= AC = \varnothing}$}
  {$\notationhighlight{F\cap J=\varnothing}$}
\probabilityrow
  {}
  {Addition for disjoint events}
  {If $E_1\cap E_2=\varnothing$ \qquad then $P(E_1\cup E_2)=P(E_1)+P(E_2)$}
  {$P(A\cup C)=\tfrac12+\tfrac12=1$}
  {$P(F\cup J)=\tfrac34+\tfrac14=1$}
\paddedcell{\centering\SymbolFont $\displaystyle
  \begin{gathered}P(E_1\cap E_2)\\{}=P(E_1)P(E_2)\end{gathered}$}&
\paddedcell{\DefinitionFont\conceptmeaning{Independent events}{%
  \meaningexplanation{The product rule at left defines independence.}\par
  \meaningexplanation{Knowing one does not change the other's probability.}}}&
\multicolumn{2}{p{\dimexpr2\ExampleWidth+2\tabcolsep\relax}@{}}{\paddedcell{\ExampleFont\raggedright
  \textbf{Coin example: independence}\par
  \begingroup
  \begin{tabularx}{\linewidth}{@{}>{\collectcell\contextcell}l<{\endcollectcell}>{\collectcell\subrowcell}X<{\endcollectcell}@{}}
  \textbf{Symbol}&\textbf{Meaning}\\
  $E_1=\{\mathrm{HH},\mathrm{HT}\}$&First toss is heads\\
  $E_2=\{\mathrm{HH},\mathrm{TH}\}$&Second toss is heads\\
  $E_1\cap E_2=\{\mathrm{HH}\}$&Both tosses are heads\\
  $P(E_1)=P(E_2)=\tfrac12$&Each event has two of four outcomes\\
  $P(E_1\cap E_2)=P(E_1)P(E_2)=\tfrac12\cdot\tfrac12=\tfrac14$ \hspace{3mm} &One of four outcomes
  \end{tabularx}\endgroup
}}\\\hline
\end{tabular}\par
\endgroup
{\color{PosterInk}\hrule height 0.5pt}\vspace{0.7mm}
\endgroup
\endgroup
\newpage
% ==========================================================================
% PAGE B-1: SEQUENCES AND SUMMATION
% ==========================================================================
\selectposterpalette{B}
\posterheader{B-1\quad SEQUENCES AND SUMMATION}
\begingroup
\compacttables
% --- Sequences and subscripts ---
\sectionheading{Sequences and subscripts}
\renewcommand{\ContextHeading}{Reading guide}
\examplecontexts{A \textbf{sequence} is an ordered list of indexed terms. Parentheses show sequences; braces show sets. Here all indices start at $i=1$, and $n\in\mathbb N$, $n\geq1$. Read each row across to compare the three sequences.}
\begingroup
\setlength{\tabcolsep}{0.8mm}
\setlength{\RowPadding}{0.95mm}
\renewcommand{\arraystretch}{1}
\arrayrulecolor{PosterRule}
% Local row helpers serve sequences and partial sums.
% Concept-only rows leave Symbol empty rather than introduce new notation.
% Coral marks isolate newly introduced symbols, indices, and bounds.
% Carry each row's role through every example, including substituted indices and values.

\newcommand{\sequencehead}{%
  \paddedcell{\centering\TableHeadFont Symbol}&%
  \paddedcell{\TableHeadFont Meaning}&%
  \paddedcell{\TableHeadFont Sequence $a$}&%
  \paddedcell{\TableHeadFont Sequence $b$}&%
  \paddedcell{\TableHeadFont Sequence $c$}\\\hline}
\begin{tabularx}{\linewidth}{@{}p{.10\linewidth}p{.18\linewidth}p{.225\linewidth}p{.225\linewidth}X@{}}
\sequencehead
\comparisonrow
  {k_{\notationhighlight{i}}=f(\notationhighlight{i})}
  {Rule and first terms}
  {Substitute the index.}
  {\examplemath[l]{a_{\notationhighlight{i}}=2\cdot\notationhighlight{i}\\(2\cdot\notationhighlight{1},2\cdot\notationhighlight{2},2\cdot\notationhighlight{3},2\cdot\notationhighlight{4},\ldots)\\=(2,4,6,8,\ldots)}}
  {\examplemath[l]{b_{\notationhighlight{i}}=\notationhighlight{i}^2\\((\notationhighlight{1})^2,(\notationhighlight{2})^2,(\notationhighlight{3})^2,(\notationhighlight{4})^2,\ldots)\\=(1,4,9,16,\ldots)}}
  {\examplemath[l]{c_{\notationhighlight{i}}=2^{\notationhighlight{i}-1}\\(2^{\notationhighlight{1}-1},2^{\notationhighlight{2}-1},2^{\notationhighlight{3}-1},2^{\notationhighlight{4}-1},\ldots)\\=(1,2,4,8,\ldots)}}
\comparisonrow
  {k_{\notationhighlight{i}}}
  {Indexed term}
  {Subscript: position.}
  {\examplemath[l]{a_{\notationhighlight{3}}=2(\notationhighlight{3})=6\\a_{\notationhighlight{4}}=2(\notationhighlight{4})=8}}
  {\examplemath[l]{b_{\notationhighlight{3}}=(\notationhighlight{3})^2=9\\b_{\notationhighlight{5}}=(\notationhighlight{5})^2=25}}
  {\examplemath[l]{c_{\notationhighlight{3}}=2^{\notationhighlight{3}-1}=4\\c_{\notationhighlight{6}}=2^{\notationhighlight{6}-1}=32}}
\comparisonrow
  {k_{\notationhighlight{1}},\ k_{\notationhighlight{n}}}
  {First term; $n$th term}
  {Term at $1$ or $n$.}
  {$a_{\notationhighlight{1}}=2$ \par $a_{\notationhighlight{n}}=2\notationhighlight{n}$}
  {$b_{\notationhighlight{1}}=1$ \par $b_{\notationhighlight{n}}=\notationhighlight{n}^2$}
  {$c_{\notationhighlight{1}}=1$ \par $c_{\notationhighlight{n}}=2^{\notationhighlight{n}-1}$}
\comparisonrow
  {(k_{i})_{i=1}^{\notationhighlight{4}}}
  {Finite sequence}
  {Last index fixes length.}
  {\examplemath[l]{(a_{i})_{i=1}^{\notationhighlight{4}}=(2,4,6,8)\\(a_{i})_{i=1}^{\notationhighlight{3}}=(2,4,6)}}
  {\examplemath[l]{(b_{i})_{i=1}^{\notationhighlight{4}}=(1,4,9,16)\\(b_{i})_{i=1}^{\notationhighlight{2}}=(1,4)}}
  {\examplemath[l]{(c_{i})_{i=1}^{\notationhighlight{4}}=(1,2,4,8)\\(c_{i})_{i=1}^{\notationhighlight{1}}=(1)}}
\comparisonrow
  {(k_{i})_{i=1}^{\notationhighlight{\infty}}}
  {Infinite sequence, etc.}
  {No last term; $\infty$ is no index.}
  {$(a_{i})_{i=1}^{\notationhighlight{\infty}}$\par $=(2,4,6,8,\ldots)$}
  {$(b_{i})_{i=1}^{\notationhighlight{\infty}}$\par $=(1,4,9,16,\ldots)$}
  {$(c_{i})_{i=1}^{\notationhighlight{\infty}}$\par $=(1,2,4,8,\ldots)$}
\comparisonrow
  {}
  {Sequence order}
  {Order matters.}
  {$(a_{1},a_{2})\ne(a_{2},a_{1})$\par $(2,4)\ne(4,2)$}
  {$(b_{1},b_{2})\ne(b_{2},b_{1})$\par $(1,4)\ne(4,1)$}
  {$(c_{1},c_{2})\ne(c_{2},c_{1})$\par $(1,2)\ne(2,1)$}
\comparisonrow
  {}
  {Sequence repetition}
  {Repetition matters.}
  {$(2,2,4)\ne(2,4)$}
  {$(1,1,4)\ne(1,4)$}
  {$(1,1,2)\ne(1,2)$}
\end{tabularx}\par
\vspace{0.5mm}
{\FooterFont An index is not multiplication or an exponent. Sets ignore order and repetition sequences do not.\par}
\vspace{3mm}
% --- Summation and partial sums ---
\sectionheading{Summation and partial sums}
\renewcommand{\ContextHeading}{From terms to their sum}
\examplecontexts{A \textbf{series} is a sum of sequence terms. The symbol $\mathop{\notationhighlight{\mathord{\sum}}}\displaylimits$ adds the summand as its index runs through consecutive integers between the bounds, \textbf{inclusive}. A partial sum adds finitely many terms.\par
In $S_{n}^{(\notationhighlight{a})}$, $S_{m}^{(\notationhighlight{b})}$, and $S_{l}^{(\notationhighlight{c})}$, the parenthesized superscripts \textbf{label the sequence being summed; they are not powers}. The subscript $\notationhighlight{n}$, $\notationhighlight{m}$, or $\notationhighlight{l}$ counts the terms when the sum starts at $1$.}
\begingroup
\begin{tabularx}{\linewidth}{@{}p{.10\linewidth}p{.2\linewidth}p{.225\linewidth}p{.225\linewidth}X@{}}
\sequencehead
\comparisonrow
  {\notationhighlight{k_{i}} = f(i)}
  {Summand}
  {Term being added.}
  {$\notationhighlight{a_{i}}=2i$}
  {$\notationhighlight{b_{i}}=i^2$}
  {$\notationhighlight{c_{i}}=2^{i-1}$}
\comparisonrow
  {\sum_{\notationhighlight{i}=\notationhighlight{1}}^{\notationhighlight{n}}f(\notationhighlight{i})}
  {Bounds and index}
  {Lower bound $1$; upper bound $n,m,l$; local index $i$.}
  {$\sum_{\notationhighlight{i}=\notationhighlight{1}}^{\notationhighlight{n}}a_{\notationhighlight{i}}$\par\vspace{1.8mm}Start: $a_{\notationhighlight{1}}$ | End: $a_{\notationhighlight{n}}$}
  {$\sum_{\notationhighlight{i}=\notationhighlight{1}}^{\notationhighlight{m}}b_{\notationhighlight{i}}$\par\vspace{1.8mm}Start: $b_{\notationhighlight{1}}$ | End: $b_{\notationhighlight{m}}$}
  {$\sum_{\notationhighlight{i}=\notationhighlight{1}}^{\notationhighlight{l}}c_{\notationhighlight{i}}$\par\vspace{1.8mm}Start: $c_{\notationhighlight{1}}$ | End: $c_{\notationhighlight{l}}$}
\comparisonrow
  {\notationhighlight{S}_{n}}
  {Partial sum of the first $n$ terms}
  {Add through index $n$, $m$, or $l$.}
  {$\notationhighlight{S}_{n}^{(a)}=\sum_{i=1}^{n}a_{i}$\par\vspace{1.8mm}\qquad\! $=a_{1}+\cdots+a_{n}$}
  {$\notationhighlight{S}_{m}^{(b)}=\sum_{i=1}^{m}b_{i}$\par\vspace{1.8mm}\qquad\! $=b_{1}+\cdots+b_{m}$}
  {$\notationhighlight{S}_{l}^{(c)}=\sum_{i=1}^{l}c_{i}$\par\vspace{1.8mm}\qquad\! $=c_{1}+\cdots+c_{l}$}
\comparisonrow
  {\notationhighlight{a}_{4}\text{ vs }\notationhighlight{S}_{4}}
  {Fourth term versus fourth partial sum}
  {One term vs. four-term sum.}
  {$\notationhighlight{a}_{4}=\notationhighlight{8}$\par\vspace{1.8mm} $\notationhighlight{S}_{4}^{(a)}=2+4+6+8=\notationhighlight{20}$}
  {$\notationhighlight{b}_{4}=\notationhighlight{16}$\par\vspace{1.8mm} $\notationhighlight{S}_{4}^{(b)}=1+4+9+16=\notationhighlight{30}$}
  {$\notationhighlight{c}_{4}=\notationhighlight{8}$\par\vspace{1.8mm} $\notationhighlight{S}_{4}^{(c)}=1+2+4+8=\notationhighlight{15}$}
\comparisonrow
  {\notationhighlight{i},\notationhighlight{j},\notationhighlight{k}}
  {Index renaming}
  {Rename the index throughout; keep bounds.}
  {\examplemath[l]{S_{4}^{(a)}=\sum_{\notationhighlight{i}=1}^{4}a_{\notationhighlight{i}}=20\\{}\qquad\!=\sum_{\notationhighlight{j}=1}^{4}a_{\notationhighlight{j}}=20\\{}\qquad\!=\sum_{\notationhighlight{k}=1}^{4}a_{\notationhighlight{k}}=20}}
  {\examplemath[l]{S_{4}^{(b)}=\sum_{\notationhighlight{i}=1}^{4}b_{\notationhighlight{i}}=30\\{}\qquad\!=\sum_{\notationhighlight{j}=1}^{4}b_{\notationhighlight{j}}=30\\{}\qquad=\sum_{\notationhighlight{k}=1}^{4}b_{\notationhighlight{k}}=30}}
  {\examplemath[l]{S_{4}^{(c)}=\sum_{\notationhighlight{i}=1}^{4}c_{\notationhighlight{i}}=15\\{}\qquad\!=\sum_{\notationhighlight{j}=1}^{4}c_{\notationhighlight{j}}=15\\{}\qquad\!=\sum_{\notationhighlight{k}=1}^{4}c_{\notationhighlight{k}}=15}}
\end{tabularx}\par
\endgroup
\vspace{3mm}
% --- Worked summation examples ---
\sectionheading{Worked summation examples}
% Label only the first expansion. Reserve the full stack depth, including
% clearance for highlight padding, a short upward arrow, and a readable label.
\newcommand{\indexedterm}[3]{\mathord{\vtop{\offinterlineskip
  \halign{\hfil##\hfil\cr
    $\mathsurround=0pt #1$\cr
    \noalign{\kern3.5pt}$\mathsurround=0pt\scriptscriptstyle\uparrow$\cr
    \noalign{\kern1.5pt}$\mathsurround=0pt\scriptstyle\notationhighlight{#2=#3}$\cr}}}}
\setlength{\RowPadding}{1.5mm}
\renewcommand{\ExampleLineStretch}{1.3}
\renewcommand{\SymbolColumnFraction}{.23}
\renewcommand{\MeaningColumnFraction}{.35}
\paralleltables{
\notationrow
  {\displaystyle\sum_{\notationhighlight{i=2}}^{4}3}
  {\conceptmeaning{Constant sum}{\meaningexplanation{Add $3$ for $\notationhighlight{i=2,3,4}$}}}
  {\examplemath[l]{\indexedterm{3}{i}{2}+\indexedterm{3}{i}{3}+\indexedterm{3}{i}{4}=9}}
\notationrow
  {\displaystyle\sum_{\notationhighlight{k=1}}^{3}k^{\notationhighlight{2}}}
  {\conceptmeaning{Sum of squares}{\meaningexplanation{Square $k$ for\\ $\notationhighlight{k=1,2,3}$}}}
  {\examplemath[l]{\indexedterm{1^{\notationhighlight{2}}}{k}{1}+\indexedterm{2^{\notationhighlight{2}}}{k}{2}+\indexedterm{3^{\notationhighlight{2}}}{k}{3}\\[1.8mm]{}=1+4+9=14}}
\notationrow
  {\displaystyle\sum_{\notationhighlight{j=0}}^{3}2^{j}}
  {\conceptmeaning{Starting at zero}{\meaningexplanation{Powers of $2$;\\ $\notationhighlight{j=0,1,2,3}$}}}
  {\examplemath[l]{\indexedterm{2^{\notationhighlight{0}}}{j}{0}+\indexedterm{2^{1}}{j}{1}+\indexedterm{2^{2}}{j}{2}+\indexedterm{2^{3}}{j}{3}\\[1.8mm]{}=1+2+4+8=15}}
}{
\notationrow
  {\displaystyle\sum_{\notationhighlight{i=3}}^{5}i}
  {\conceptmeaning{Shifted lower bound}{\meaningexplanation{Add $i$ for $\notationhighlight{i=3,4,5}$}}}
  {\examplemath[l]{\indexedterm{\notationhighlight{3}}{i}{3}+\indexedterm{4}{i}{4}+\indexedterm{5}{i}{5}=12}}
\notationrow
  {\displaystyle\sum_{\notationhighlight{j=2}}^{4}b_{\notationhighlight{j}}}
  {\conceptmeaning{Sequence terms}{\meaningexplanation{$b_{\notationhighlight{j}}=\notationhighlight{j}^2$;\\ $\notationhighlight{j=2,3,4}$}}}
  {\examplemath[l]{\indexedterm{b_{\notationhighlight{2}}}{j}{2}+\indexedterm{b_{\notationhighlight{3}}}{j}{3}+\indexedterm{b_{\notationhighlight{4}}}{j}{4}\\[1.8mm]{}=(\notationhighlight{2})^2+(\notationhighlight{3})^2+(\notationhighlight{4})^2\\[1.8mm]{}=4+9+16=29}}
\notationrow
  {\displaystyle\sum_{\notationhighlight{k=1}}^{4}(2\notationhighlight{k}+1)}
  {\conceptmeaning{Linear summand}{\meaningexplanation{$\notationhighlight{k=1,\ldots,4}$; add $2\notationhighlight{k}+1$}}}
  {\examplemath[l]{\indexedterm{(2(\notationhighlight{1})+1)}{k}{1}+\indexedterm{(2(\notationhighlight{2})+1)}{k}{2}\\[1.8mm]{}\quad+\indexedterm{(2(\notationhighlight{3})+1)}{k}{3}+\indexedterm{(2(\notationhighlight{4})+1)}{k}{4}\\[1.8mm]{}=3+5+7+9=24}}
}
\endgroup
{\color{PosterInk}\hrule height 0.5pt}\vspace{0.7mm}
{\FooterFont Symbols acquire meaning from their mathematical context. State the starting index and label the sequence in every partial sum.\par}
\endgroup
\newpage
% ==========================================================================
% PAGE B-2: LIMITS AND CALCULUS
% ==========================================================================
\selectposterpalette{B}
\posterheader{B-2\quad LIMITS AND CALCULUS}
\begingroup
% The densest poster needs slightly smaller body/example type; symbols and
% headings retain the common sizes and no table is scaled.
\renewcommand{\DefinitionFont}{\fontsize{12.6}{14.4}\selectfont}
\renewcommand{\ExampleFont}{\fontsize{12.4}{14.4}\selectfont}
\setlength{\RowPadding}{0.20mm}
% Local column proportions reserve room for wrapped calculus examples.
\renewcommand{\SymbolColumnFraction}{.25}
\renewcommand{\MeaningColumnFraction}{.34}
% Body notationrow uses centeredcell directly; paddedcell remains a separate
% dispatch hook so headers/references keep 0.20 mm while body clearance varies.
% These definitions and all smaller calculus fonts are confined to this page.
\newcommand{\calculusrowpadding}[1]{%
  \setlength{\RowPadding}{#1}%
  \renewcommand{\paddedcell}[1]{{\setlength{\RowPadding}{0.20mm}\centeredcell{##1}}}%
}
% --- Limits and approaching values: read the numbered columns left to right ---
\sectionheading{Limits and approaching values}
\begingroup
\setlength{\tabcolsep}{2mm}
\arrayrulecolor{PosterRule}
% Three measured cells give the progression a visible order; no table scaling.
\begin{tabularx}{\linewidth}{@{}p{.25\linewidth}p{.37\linewidth}X@{}}
\paddedcell{\TableHeadFont 1. Approaching an input}&
\paddedcell{\TableHeadFont 2. Approaching an output}&
\paddedcell{\TableHeadFont 3. Distinguish special cases}\\\hline
\paddedcell{\DefinitionFont
  $x\mathrel{\notationhighlight{\mathord{\to}}}a$: inputs get arbitrarily close to $a$; this is different from $x=a$.\par\vspace{1.5mm}
  $x\mathrel{\notationhighlight{\mathord{\to}}}2$: $x$ approaches $2$; $x=2$ is exact.\par\vspace{2mm}
  \textbf{Read the two sides together:}\par
  $x\to a^{\notationhighlight{-}}$: approach from $x<a$.\par
  $x\to a^{\notationhighlight{+}}$: approach from $x>a$.\par\vspace{1.5mm}
  For example, $x\to2^{\notationhighlight{-}}$ uses $x<2$; $x\to2^{\notationhighlight{+}}$ uses $x>2$.\par\vspace{1.5mm}
  \meaningexplanation{$a^-$ and $a^+$ show the approach side, not the sign of $f(x)$.}}
&\paddedcell{\DefinitionFont
  {\centering\SymbolFont $\displaystyle
    \mathop{\notationhighlight{\mathord{\lim}}}\limits_{x\to a}f(x)=L$\par}
  Read: input $x$ approaches $a$; output $f(x)$ approaches $L\in\mathbb R$.\par\vspace{0.6mm}
  Use nearby domain inputs $x\ne a$; $f(a)$ may differ from $L$ or be undefined.\par\vspace{0.6mm}
  \textbf{Both sides must agree:}\par
  $\lim_{x\to a^-}f(x)=L=\lim_{x\to a^+}f(x)$\par
  means the two-sided limit is $L$.\par\vspace{0.6mm}
  \ExampleFont
  $\mathop{\notationhighlight{\mathord{\lim}}}_{x\to2^-}x^2=4=\mathop{\notationhighlight{\mathord{\lim}}}_{x\to2^+}x^2$, so $\mathop{\notationhighlight{\mathord{\lim}}}_{x\to2}x^2=4$\par\vspace{0.6mm}
  $\lim_{x\to1}\frac{x^2-1}{x-1}=2$: the quotient equals $x+1$ for $x\ne1$, but is undefined at $x=1$.}
&\paddedcell{\DefinitionFont
  \textbf{Unequal one-sided limits:}\par
  For $g(x)=|x|/x$, $x\ne0$,\par
  $\lim_{x\to0^-}g(x)=-1\ne1=\lim_{x\to0^+}g(x)$\par
  The two-sided limit does not exist ($\notationhighlight{\text{DNE}}$).\par\vspace{0.6mm}
  \textbf{Inputs without bound:}\par
  $x\to\notationhighlight{+\infty}$ or $x\to\notationhighlight{-\infty}$: no finite target $a$.\par
  $\lim_{x\to+\infty}\frac1x=0=\lim_{x\to-\infty}\frac1x$\par\vspace{0.6mm}
  \textbf{Outputs without bound:}\par
  $f(x)\to\notationhighlight{+\infty}$ or $f(x)\to\notationhighlight{-\infty}$: infinite-limit statements, not finite real limits.\par
  \ExampleFont $\lim_{x\to a}\frac1{(x-a)^2}=+\infty$;\par
  $\lim_{x\to0^+}\frac1x=+\infty$, $\lim_{x\to0^-}\frac1x=-\infty$}
\\\hline
\end{tabularx}\par
\endgroup
\vspace{2mm}
% --- Change and derivatives ---
\sectionheading{Change and derivatives}
\renewcommand{\ContextHeading}{A derivative is a limit}
\examplecontexts{Instantaneous change:
$\displaystyle f^{\notationhighlight{\prime}}(a)=\lim_{h\to0}\frac{f(a+h)-f(a)}h$, $h\ne0$; the two-sided limit must be finite.
For $f(x)=x^2$: $f'(a)=\lim_{h\to0}(2a+h)=2a$}
\noindent\begin{minipage}[t]{.49\linewidth}\vspace{0pt}
\notationtable{
\notationrow
  {\begin{gathered}\notationhighlight{\Delta} x,\\\notationhighlight{\Delta} y\end{gathered}}
  {\conceptmeaning{Finite changes}{\meaningexplanation{$\notationhighlight{\Delta} x=x_2-x_1$; $\notationhighlight{\Delta} y=f(x_2)-f(x_1)$}}}
  {$f(x)=x^2$, from $x=1$ to $x=3$:\par $\notationhighlight{\Delta} x=2$, $\notationhighlight{\Delta} y=9-1=8$}
\notationrow
  {\frac{\notationhighlight{\Delta} y}{\notationhighlight{\Delta} x}}
  {\conceptmeaning{Average rate of change}{\meaningexplanation{Output change per input change; $\notationhighlight{\Delta} x\ne0$.}}}
  {$f(x)=x^2$, from $1$ to $3$:\par $\notationhighlight{\Delta} y/\notationhighlight{\Delta} x=8/2=4$}
\notationrow
  {\begin{gathered}f^{\notationhighlight{\prime}}(x),\quad y^{\notationhighlight{\prime}},\\\frac{\notationhighlight{d}y}{\notationhighlight{d}x}\end{gathered}}
  {\conceptmeaning{First derivative}{\meaningexplanation{Instantaneous rate for $y=f(x)$.}}}
  {$y=f(x)=x^2$:\par $\begin{gathered}f^{\notationhighlight{\prime}}(x)=y^{\notationhighlight{\prime}}=\dfrac{\notationhighlight{d}y}{\notationhighlight{d}x}=2x\end{gathered}$}
}
\end{minipage}\hfill
\begin{minipage}[t]{.49\linewidth}\vspace{0pt}
\calculusrowpadding{1mm}
\notationtable{
\notationrow
  {\frac{\notationhighlight{d}}{\notationhighlight{d}x},\quad f^{\notationhighlight{\prime}}(a)}
  {\conceptmeaning{Differentiation; value at a point}{\meaningexplanation{$\notationhighlight{d}/\notationhighlight{d}x$: operator; $f^{\notationhighlight{\prime}}(a)$: derivative value at $a$.}}}
  {$\dfrac{\notationhighlight{d}}{\notationhighlight{d}x}(x^2)=2x$\par $f^{\notationhighlight{\prime}}(3)=2(3)=6$}
\notationrow
  {\begin{gathered}f^{\notationhighlight{\prime\prime}}(x),\\\frac{d^{\notationhighlight{2}}y}{dx^{\notationhighlight{2}}}\end{gathered}}
  {\conceptmeaning{Second derivative}{\meaningexplanation{Derivative of the first derivative.}}}
  {$y=f(x)=x^2$:\par $\begin{gathered}f^{\notationhighlight{\prime\prime}}(x)=\dfrac{d^{\notationhighlight{2}}y}{dx^{\notationhighlight{2}}}=2\end{gathered}$}
}
\end{minipage}\par\vspace{2mm}
% --- Integrals and evaluation ---
\sectionheading{Integrals and evaluation}
\noindent\begin{minipage}[t]{.49\linewidth}\vspace{0pt}
\notationtable{
\notationrow
  {\notationintegral{\int} f(x)\,dx}
  {\conceptmeaning{Indefinite integral}{\meaningexplanation{Antiderivatives $F(x)+C$; $F'=f$, $C\in\mathbb R$.}}}
  {$\displaystyle\notationintegral{\int} 2x\,dx=x^2+C$\par A family of functions, one for each $C$.}
\notationrow
  {\notationintegral{\int_a^b} f(x)\,dx}
  {\conceptmeaning{Definite integral}{\meaningexplanation{Signed accumulation from $a$ to $b$; integrand $f(x)$.}}}
  {$\displaystyle\notationintegral{\int_1^2} 2x\,dx=3$\par A numerical value; signed accumulation is not always geometric area.}
\notationrow
  {F(x)\notationhighlight{\big|}_a^b}
  {\conceptmeaning{Evaluation at the bounds}{\meaningexplanation{$F(b)-F(a)$. Equals $\int_a^b f(x)\,dx$ if $f$ is continuous on $[a,b]$ and $F'=f$.}}}
  {\examplemath[l]{\displaystyle\int_1^2 2x\,dx=x^2\notationhighlight{\big|}_1^2\\{}=2^2-1^2=3}}
}
\end{minipage}\hfill
\begin{minipage}[t]{.49\linewidth}\vspace{0pt}
\calculusrowpadding{1mm}
\notationtable{
\notationrow
  {\begin{gathered}\notationhighlight{\Delta x}=\frac{b-a}{n}\\
  \notationhighlight{x_i}=a+i\notationhighlight{\Delta x}\\
  \notationhighlight{R_n}=\\\sum_{i=1}^n f(\notationhighlight{x_i})\notationhighlight{\Delta x}\end{gathered}}
  {\conceptmeaning{Riemann sum}{\meaningexplanation{Continuous $f$ on $[a,b]$; $n$ equal widths $\notationhighlight{\Delta x}$, right endpoints $\notationhighlight{x_i}$.}\par
  $\begin{gathered}\displaystyle\int_a^b f(x)\,dx\\{}=\lim\limits_{n\to\infty}\notationhighlight{R_n}\end{gathered}$}}
  {For $f(x)=2x$ on $[1,2]$:\par
   \examplemath[l]{\notationhighlight{\Delta x}=\notationhighlight{1/n},\\\notationhighlight{x_i}=\notationhighlight{1+i/n}\\
   \notationhighlight{R_n}=\sum_{i=1}^n2\bigl(\notationhighlight{1+i/n}\bigr)\bigl(\notationhighlight{1/n}\bigr)\\
   {}=\notationhighlight{3+1/n}\to\notationhighlight{3}}}
\notationrow
  {\notationhighlight{d}x}
  {\conceptmeaning{Integration variable}{\meaningexplanation{Integrate with respect to $x$.}}}
  {In $\int_1^2 2x\,\notationhighlight{d}x$: bounds $1$, $2$;\par integrand $2x$; integration variable $x$.}
}
\end{minipage}\par

\par\vspace{2mm}
% --- Three complementary calculus pictures: definitions above, pictures below ---
\sectionheading{Calculus in pictures}
\noindent\begin{minipage}[t]{.325\linewidth}\vspace{0pt}
{\TableHeadFont Derivative: secant to tangent\par}
\ExampleFont $f(x)=x^2$, $a=1$; pictured $h=0.6$.\par
% Derivative graph: Q=(1+h,f(1+h)); secant y=(2+h)x-(1+h).
% All coordinates and both lines are evaluated from f(x)=x^2 at a=1.
\begin{tikzpicture}[font=\fontsize{10.8}{12}\selectfont,
  every node/.style={inner sep=1pt},line cap=round]
  \path[use as bounding box] (0,0) rectangle (91mm,36mm);
  \def\secantStep{0.6}
  \begin{scope}[shift={(8mm,6mm)},x=36mm,y=7.5mm]
    \coordinate (derivativeP) at (1,1);
    \coordinate (derivativeQ) at ({1+\secantStep},{(1+\secantStep)^2});
    \draw[->,line width=0.70pt] (0,0)--(2,0) node[right] {$x$};
    \draw[->,line width=0.70pt] (0,0)--(0,3.6) node[above] {$y$};
    \draw[PosterAccent,line width=1.70pt,domain=0:1.85,samples=50]
      plot (\x,{\x*\x});
    \draw[PosterInk,dashed,line width=1.4pt,domain=.25:2]
      plot (\x,{2*\x-1});
    \draw[PosterInk!60,line width=1.2pt,domain=.5:1.95]
      plot (\x,{(2+\secantStep)*\x-(1+\secantStep)});
    \draw[PosterInk!60,<->,line width=0.70pt]
      (1,1)--({1+\secantStep},1) node[midway,below=1mm] {$\Delta x = h$};
    \draw[PosterInk!60,<->,line width=0.70pt]
      ({1+\secantStep},1)--(derivativeQ);
    \fill (derivativeP) circle[radius=1.2pt] (derivativeQ) circle[radius=1.2pt];
    \node (derivativePLabel) at (.55,2.25) {$P=(1,1)$};
    \draw[PosterInk!60,line width=0.6pt,shorten <=1mm,shorten >=1mm]
      (derivativePLabel.south)--(derivativeP);
    \node (derivativeQLabel) at (1.15,3.6) {$Q$};
    \draw[PosterInk!60,line width=0.6pt,shorten <=1mm,shorten >=1mm]
      (derivativeQLabel.south)--(derivativeQ);

    \node[anchor=west,text=PosterInk!60] at (1.72,1.7) {$\Delta y$};
    % The existing movement arrow below Delta x = h indicates h -> 0+.
    \draw[->,PosterAccent,line width=1.6pt] (1.5,.2)--(1.18,.2);
    \node[below] at (2,0) {$2$};
    \node[below left] at (0,0) {$0$};
  \end{scope}
\end{tikzpicture}\par
$Q=(1+h,(1+h)^2)$; $\Delta y=f(1+h)-f(1)$\par
Solid secant slope $\frac{\Delta y}h=2+h\to2$ as $h\to0$, $h\ne0$, from either side. $Q\to P$; dashed tangent: $y=2x-1$, $f'(1)=2$.
\end{minipage}\hfill
\begin{minipage}[t]{.31\linewidth}\vspace{0pt}
{\TableHeadFont Indefinite integral: a family\par}
\ExampleFont $\displaystyle\int2x\,dx=x^2+C$\par
% Antiderivative graph: three translations, identified by their constants.
\begin{tikzpicture}[font=\fontsize{10.8}{12}\selectfont,
  every node/.style={inner sep=1pt},line cap=round,
  densely dotted/.style={dash pattern=on .8pt off 1pt}]
  \path[use as bounding box] (0,0) rectangle (87mm,36mm);
  \begin{scope}[shift={(36mm,10mm)},x=17mm,y=6mm]
    \draw[->,line width=0.70pt] (-1.8,0)--(2.55,0) node[right] {$x$};
    \draw[->,line width=0.70pt] (0,-1.4)--(0,3.5) node[above] {$y$};
    \foreach \constant/\linestyle in {-1/densely dotted,0/solid,1/dashed}{
      \draw[PosterAccent,\linestyle,line width=1.6pt,domain=-1.4:1.4,samples=45]
        plot (\x,{\x*\x+\constant});
      \node[anchor=west] at (1.45,{1.4^2+\constant}) {$C=\constant$};
    }
    \node[below left] at (0,0) {$0$};
  \end{scope}
\end{tikzpicture}\par
All curves have derivative $2x$: $\frac{d}{dx}(x^2+C)=2x$. The indefinite integral is a \textbf{family of functions}; changing $C$ shifts a curve vertically.
\end{minipage}\hfill
\begin{minipage}[t]{.325\linewidth}\vspace{0pt}
{\TableHeadFont Definite integral: accumulation\par}
\ExampleFont $\displaystyle\int_1^2 2x\,dx=3$\par
% Definite-integral graph: exact trapezoid under y=2x, 1<=x<=2.
\begin{tikzpicture}[font=\fontsize{10.8}{12}\selectfont,
  every node/.style={inner sep=1pt},line cap=round]
  \path[use as bounding box] (0,0) rectangle (91mm,36mm);
  \begin{scope}[shift={(17mm,6mm)},x=23mm,y=6mm]
    \fill[PosterAccent!18] (1,0)--(2,0)--(2,{2*2})--(1,{2*1})--cycle;
    \draw[PosterRule,line width=0.8pt] (1,{2*1})--(1,0)--(2,0)--(2,{2*2});
    \draw[->,line width=0.70pt] (0,0)--(2.7,0) node[right] {$x$};
    \draw[->,line width=0.70pt] (0,0)--(0,4.6) node[above] {$y$};
    \draw[PosterAccent,line width=1.70pt,domain=0:2.2]
      plot (\x,{2*\x});
    \node[above left] at (1.9,3.8) {$y=2x$};
    \node at (1.5,1.1) {$3$};
    \node[below=1mm] at (0,0) {$0$};
    \node[below=1mm] at (1,0) {$1$};
    \node[below=1mm] at (2,0) {$2$};
    \node[left] at (0,4) {$4$};
    \draw[PosterRule,dash pattern=on .3pt off 1pt,line width=0.6pt] (0,4)--(2,4);
  \end{scope}
\end{tikzpicture}\par
On $[1,2]$, $2x\ge0$: the integral equals shaded area $\frac12(2+4)(2-1)=3$. In general, area below the axis counts negatively in \textbf{signed accumulation}.
\end{minipage}\par\vspace{2mm}
% --- Relevant examples ---
\sectionheading{Relevant examples}
\begingroup
\setlength{\tabcolsep}{1mm}
\setlength{\RowPadding}{0.45mm}
\renewcommand{\arraystretch}{1}
\arrayrulecolor{PosterRule}
\newcommand{\calculusrow}[3]{%
  \paddedcell{\SymbolFont $\displaystyle #1$}&%
  \paddedcell{\SymbolFont $\displaystyle #2$}&%
  \paddedcell{\SymbolFont $\displaystyle #3$}\\\hline}
\noindent\begin{minipage}[t]{.52\linewidth}\vspace{0pt}
\begin{tabularx}{\linewidth}{@{}p{.30\linewidth}p{.26\linewidth}X@{}}
\paddedcell{\TableHeadFont Function $f(x)$}&
\paddedcell{\TableHeadFont Derivative $f'(x)$}&
\paddedcell{\TableHeadFont Indefinite integral\par $\int f(x)\,dx$}\\\hline
\calculusrow
  {e^x}
  {e^x}
  {e^x+C}
\calculusrow
  {\ln x\quad(x>0)}
  {\frac1x}
  {x\ln x-x+C}
\calculusrow
  {P(x)=\sum_{k=0}^{n}a_kx^k}
  {\sum_{k=1}^{n}k a_kx^{k-1}}
  {\sum_{k=0}^{n}\frac{a_k}{k+1}x^{k+1}+C}
\end{tabularx}
\end{minipage}\hfill
\begin{minipage}[t]{.46\linewidth}\vspace{0pt}
\begin{tabularx}{\linewidth}{@{}p{.28\linewidth}p{.29\linewidth}X@{}}
\paddedcell{\TableHeadFont Function $f(x)$\par $x$ in radians}&
\paddedcell{\TableHeadFont Derivative $f'(x)$}&
\paddedcell{\TableHeadFont Indefinite integral\par $\int f(x)\,dx$}\\\hline
\calculusrow
  {\sin x}
  {\cos x}
  {-\cos x+C}
\calculusrow
  {\cos x}
  {-\sin x}
  {\sin x+C}
\calculusrow
  {\begin{gathered}\tan x\\(\cos x\ne0)\end{gathered}}
  {\begin{gathered}\sec^2 x=\\\frac1{\cos^2 x}\end{gathered}}
  {-\ln|\cos x|+C}
\end{tabularx}
\end{minipage}\par\vspace{0.8mm}
\endgroup
{\color{PosterInk}\hrule height 0.5pt}\vspace{0.7mm}
{\FooterFont Limits describe approaching behavior; derivatives describe local change; definite integrals describe signed accumulation.\par}
\endgroup
\end{document}
